\exercisegroup

\begin{enumerate}
\def\labelenumi{\arabic{enumi})}
\tightlist
\item
  Let E and F be two infinite dimensional Hilbert spaces satisfying the first axiom of countability, \((a_{n})\) an orthonormal basis of E, and \((b_{n})\) an orthonormal basis of F.
\end{enumerate}

\begin{enumerate}
\def\labelenumi{\alph{enumi})}
\tightlist
\item
  Let u be a continuous linear mapping from E into F; let \(u(a_{n}) = \sum_{m} \alpha_{mn} b_{m}\) . Show that
\[
\sum_ {n} | \alpha_ {m n} | ^ {2} \leqslant \| u \| ^ {2}
\]

\[
\sum_ {m} | \alpha_ {m n} | ^ {2} \leqslant \| u \| ^ {2}
\]
for every \(m\) and \(n\) .

\begin{enumerate}
\def\labelenumi{\alph{enumi})}
\setcounter{enumi}{1}
\tightlist
\item
  Give an example of a double sequence \((\alpha_{mn})\) such that \(\sum_{m} |\alpha_{mn}|^{2} \leqslant 1\) for all \(n\) and \(\sum_{n} |\alpha_{mn}|^{2} \leqslant 1\) for all \(m\) , but such that there does not exist any continuous linear mapping \(u\) from E into F such that \(\langle u(a_n)|b_m\rangle = \alpha_{mn}\) for every pair of integers \((m, n)\) . (Show that if I \(\subset\) N is a set of \(p\) integers, and if V \(_p\) (resp. W \(_p\) ) is the subspace of E (resp. F) generated by the \(a_n\) (resp. \(b_n\) ) such that \(n \in I\) , then there exists a linear mapping \(u_p\) from V \(_p\) onto W \(_p\) such that \(\langle u_p(a_n)|b_m\rangle = \frac{1}{\sqrt{p}}\) for \(m \in I\) and \(n \in I\) , and that \(\|u_p\| \geqslant \sqrt{p}\) .)
\end{enumerate}

\begin{enumerate}
\def\labelenumi{\arabic{enumi})}
\setcounter{enumi}{1}
\tightlist
\item
  Let \(A = (\alpha_{mn})_{(m,n)\in \mathbf{N}\times \mathbf{N}}\) be a double sequence of complex numbers, which we also call an infinite matrix. For every point \(x = (x_{n})\) of the direct sum space \(\mathbf{C}^{(\mathbf{N})}\) , the sums \(y_{m} = \sum_{n}\alpha_{mn}x_{n}\) are defined; let \(A.x\) be the point \((y_{m})\) of the product space \(\mathbf{C}^{\mathbf{N}}\) , then \(x\mapsto A.x\) is a linear mapping from \(\mathbf{C}^{(\mathbf{N})}\) into \(\mathbf{C}^{\mathbf{N}}\) , and every linear mapping from \(\mathbf{C}^{(\mathbf{N})}\) into \(\mathbf{C}^{\mathbf{N}}\) is of this form. Let \(E_{n}\) denote the subspace of \(\mathbf{C}^{(\mathbf{N})}\) generated by the first \(n\) vectors of the canonical basis, \(P_{n}\) the canonical projection from \(\mathbf{C}^{\mathbf{N}}\) onto \(E_{n}\) ; when \(E_{n}\) is assigned the norm induced by that of the space \(\ell_{\mathbf{C}}^{2}(\mathbf{N})\) , \(\| u\|\) denotes the norm of the linear mapping \(u\) from the finite dimensional hilbertian space \(E_{n}\) into itself.
\end{enumerate}

\begin{enumerate}
\def\labelenumi{\alph{enumi})}
\item
  In order that the image of \(\mathbf{C}^{(\mathbf{N})}\) under the mapping \(x\mapsto A.x\) be contained in \(\ell_{\mathbf{C}}^2 (\mathbf{N})\) and that this mapping extend to a continuous linear mapping from \(\ell_{\mathbf{C}}^2 (\mathbf{N})\) into itself, it is necessary and sufficient that the norms \(\| P_nAP_n\|\) are bounded. This implies that the rows and the columns of \(A\) belong to \(\ell_{\mathbf{C}}^2 (\mathbf{N})\) (exerc. 1).
\item
  Let \(A^*\) denote the infinite matrix \((\alpha_{mn}')\) , where \(\alpha_{mn}' = \overline{\alpha}_{nm}\) . If the columns of \(A\) belong to \(\ell_{\mathbf{C}}^2(\mathbf{N})\) (in other words, if \(x \mapsto A.x\) maps \(\mathbf{C}^{(\mathbf{N})}\) into \(\ell_{\mathbf{C}}^2(\mathbf{N})\) ), then the series \(\beta_{mn} = \sum_{p} \overline{\alpha}_{pm} \alpha_{pn}\) are
\end{enumerate}

absolutely convergent, and we put \(A^{*}A = (\beta_{mn})\) . Show that for \(x \mapsto A.x\) to extend to a continuous linear mapping \(u\) from \(\ell_{\mathbf{C}}^{2}(\mathbf{N})\) into itself, it is necessary and sufficient that the norms \(\| P_n(A^*A)P_n\|\) are bounded (we have \(\langle P_n(A^*A)P_n.x|x\rangle = \| AP_n .x\|^2\) for all \(x\in E_n\) ). Then \(x\mapsto A^{*}A.x\) extends to a positive hermitian mapping \(u^{*}u\) from \(\ell_{\mathbf{C}}^{2}(\mathbf{N})\) into itself.

\begin{enumerate}
\def\labelenumi{\alph{enumi})}
\setcounter{enumi}{2}
\item
  For two infinite matrices \(X = (\xi_{mn})\) , \(Y = (\eta_{mn})\) , we say that the product \(XY\) is defined if the series \(\zeta_{mn} = \sum_{p} \xi_{mp} \eta_{pn}\) are absolutely convergent, and then we put \(XY = (\zeta_{mn})\) . We say that a power \(X^k\) ( \(k\) integer \(> 1\) ) is defined if \(X^{k-1}\) and \(X^{k-1}X\) are defined and then we put \(X^k = X^{k-1}X\) ; in this case, \(X^pX^q = X^k\) for every pair of integers \(p, q\) such that \(p + q = k\) . If \(A\) is an infinite matrix whose columns are in \(\ell_{\mathbf{C}}^2(\mathbf{N})\) and if the product \((A^*A)^2\) is defined, show that for all \(x \in E_n\) , we have \(\langle (P_nA^*AP_n)^2.x|x\rangle \leqslant \langle P_n(A^*A)^2P_n.x|x\rangle\) , and deduce that \(\| P_nA^*AP_n\|^2 \leqslant \| P_n(A^*A)^2P_n\|\) .
\item
  In order that an infinite matrix A be such that the image of \(\mathbf{C}^{(\mathbf{N})}\) under \(x\mapsto A.x\) is contained in \(\ell_{\mathbf{C}}^2 (\mathbf{N})\) and that \(x\mapsto A.x\) extends to a continuous linear mapping from \(\ell_{\mathbf{C}}^2 (\mathbf{N})\) into itself, it is necessary and sufficient that the following three conditions are satisfied:
\end{enumerate}

\begin{enumerate}
\def\labelenumi{(\roman{enumi})}
\item
  the rows and columns of \(A\) are in \(\ell_{\mathbf{C}}^{2}(\mathbf{N})\) ;
\item
  the powers \((A^{*}A)^{k}\) are defined for every integer \(k > 1\) ;
\item
  we have
\end{enumerate}

\[
\sup _ {n} \left(\sup _ {m} | ((A ^ {*} A) _ {m m} ^ {n}) ^ {1 / n} |\right) <   + \infty
\]

where \((A^{*}A)_{mm}^{n}\) denotes the term with indices \((m,m)\) of the matrix \((A^{*}A)^{n}\) . (Observe that if \(C\) is the matrix, with respect to the canonical basis of \(\mathrm{E}_n\) , of a positive hermitian endomorphism of \(\mathrm{E}_n\) , then \(\| C\| \leqslant n\sup_{1\leqslant i\leqslant n}|C_{ii}|\) , by considering the trace of \(C\) and by diagonalizing \(C\) . Using the inequality proved in \(c\) ), show that

\[
\left\| P _ {n} A ^ {*} A P _ {n} \right\| \leqslant n ^ {2 - k} \sup _ {1 \leqslant i \leqslant n} \left| ((A ^ {*} A) _ {i i} ^ {2 k}) \right| ^ {2 - k}
\]

for every integer \(k > 1\) , if conditions (i), (ii) and (iii) are satisfied.)

¶ 3) Let \((a_{ij})_{(i,j)\in\mathbb{I}\times\mathbb{I}}\) be a countably infinite double family of complex numbers. Assume that there exist two numbers \(\beta > 0\) , \(\gamma > 0\) , and a family \((p_i)_{i\in\mathbb{I}}\) of numbers > 0 satisfying the relations

\[
\sum_ {i} p _ {i} | a _ {i j} | \leqslant \beta p _ {j}, \quad \sum_ {j} | a _ {i j} | p _ {j} \leqslant \gamma p _ {i}\tag{*}
\]

for all \(i, j\) in \(\mathbb{I}\) .
a) Show that there exists a continuous endomorphism \(u\) from \(\ell_{\mathbf{C}}^2 (\mathrm{I})\) , with norm \(\leqslant (\beta \gamma)^{1 / 2}\) , such that for all \(x = (x_{i})_{i\in \mathbf{I}}\) in \(\ell_{\mathbf{C}}^2 (\mathrm{I})\) , we have \(u(x) = y\) , where \(y = (y_{i})\) is given by \(y_{i} = \sum_{j}a_{ij}x_{j}\) (for \(x = (x_{i})\) and \(y = (y_{i})\) in \(\mathbf{C}^{(\mathbb{I})}\) , put \(v_{ij} = |x_i| (p_j|a_{ij}| / p_i)^{1 / 2}\) , \(w_{ij} = |y_j|(p_i|a_{ij}| / p_j)^{1 / 2}\) , and find a bound for \(\sum_{i,j}v_{ij}w_{ij}\) ).

\begin{enumerate}
\def\labelenumi{\alph{enumi})}
\setcounter{enumi}{1}
\tightlist
\item
  For I, take the set of all integers \(\geqslant 1\) , and let \(a_{ij} = (i + j)^{-1}\) . Show that the conditions (*) are satisfied with \(p_i = i^{-1/2}\) and \(\beta = \gamma = \pi\) (these being the best possible constants) (compare the series in (*) with an integral). In an analogous manner, treat the case where I = N and \(a_{ij} = (i + j + 1)^{-1}\) (« Hilbert's matrix »).
\end{enumerate}

* c) Let \(\mathcal{H} = \mathrm{H}^2 (\mathbf{D})\) be the Hardy space, consisting of all functions \(f(z) = \sum_{n=0}^{\infty} a_n z^n\) , analytic in the open disc \(\mathbf{D}:|z| < 1\) and such that \(\| f\|^2 = \sum_n |a_n|^2 < +\infty\) ; then \(\| f\|\) is a norm on \(\mathcal{H}\) , for which \(\mathcal{H}\) is isomorphic to \(\ell_{\mathbf{C}}^2 (\mathbf{N})\) . Given two functions \(f, g\) in \(\mathcal{H}\) , show that the function \(t \mapsto f(t) g(t)\) of the real variable \(t\) is integrable on \([0,1]\) with respect to the Lebesgue measure and that the formula \(\mathbf{B}(f,g) = \int_0^1 f(t) g(t) dt\) defines a continuous bilinear form on \(\mathcal{H} \times \mathcal{H}\) (consider \(\mathbf{B}(f,f) = \int_0^1 f(t)^2 dt\) for a function \(f \in \mathcal{H}\) of the form \(f(z) = \sum_{n=0}^{\mathbf{N}} a_n z^n\) with \(a_n \geqslant 0\) for \(0 \leqslant n \leqslant \mathbf{N}\) ; use Cauchy's theorem to establish the relation

\[
\int_ {- 1} ^ {1} f (t) ^ {2} d t = - i \int_ {0} ^ {\pi} f (e ^ {i \theta}) ^ {2} e ^ {i \theta} d \theta
\]

from which we get \(\mathbf{B}(f, f) \leqslant \frac{1}{2} \int_{-\pi}^{\pi} |f(e^{i\theta})|^2 d\theta = \pi \| f\|^2\) . In this way, get the result of \(b\) , according to which the Hilbert's matrix defines an endomorphism of norm \(\leqslant \pi\) of the hilbertian space \(\ell_{\mathbf{C}}^2(\mathbf{N})\) . *

\begin{enumerate}
\def\labelenumi{\arabic{enumi})}
\setcounter{enumi}{3}
\tightlist
\item
  Let \(\mathbf{E}\) be a complex hilbertian space of finite dimension \(d\) .
\end{enumerate}

\begin{enumerate}
\def\labelenumi{\alph{enumi})}
\item
  For every \(u \geqslant 0\) in \(\mathcal{L}(\mathrm{E})\) (V, p.~45), prove that there exists a unique \(v \geqslant 0\) in \(\mathcal{L}(\mathrm{E})\) such that \(u = v^2\) (diagonalize \(u\) ); we write \(v = u^{1/2}\) .
\item
  For every \(u \in \mathcal{L}(\mathrm{E})\) , put \(\mathrm{abs}(u) = (u^{*}u)^{1/2}\) . Show that u and \(\mathrm{abs}(u)\) have the same norm and that we have \(\mathrm{abs}(\symbf{\Lambda}^{n}(u)) = \symbf{\Lambda}^{n}(\mathrm{abs}(u))\) for every integer \(n \leqslant d\) .
\item
  Let \(s_{1}(u) \geqslant s_{2}(u) \geqslant \ldots \geqslant s_{d}(u) \geqslant 0\) be the sequence of eigenvalues of \(\mathrm{abs}(u)\) counted with their order of multiplicity. Show that \(\|u\| = s_{1}(u)\) and that for every integer \(n \leqslant d\) , \(\| \symbf{\Lambda}^{n}(u)\| = s_{1}(u) s_{2}(u) \ldots s_{n}(u)\) . In order that \(\| \symbf{\Lambda}^{n}(u)\| = \| u\|^{n}\) for all n such that \(1 \leqslant n \leqslant d\) , it is necessary and sufficient that \(u^{*}u\) is a homothety, in other words, that u is a scalar multiple of a unitary operator.
\end{enumerate}

\begin{enumerate}
\def\labelenumi{\arabic{enumi})}
\setcounter{enumi}{4}
\tightlist
\item
  Let E, F be two Hilbert spaces, and u a continuous linear mapping from E into F. Let \(\ell(u)\) denote the set of all \(x \in E\) such that \(\|u(x)\| = \|u\|. \|x\|\) .
\end{enumerate}

\begin{enumerate}
\def\labelenumi{\alph{enumi})}
\item
  Show that \(\ell(u)\) is the closed vector subspace of E which is the kernel of \(u^{*}u - \| u\|^{2}1_{\mathrm{E}}\) and is orthogonal to the kernel of \(u\) .
\item
  Show that the restriction of \(u\) to \(\ell(u)\) is a bijection from \(\ell(u)\) onto \(\ell(u^*)\) , whose inverse bijection is the restriction of \(\|u\|^{-2}.u^*\) to \(\ell(u^*)\) ; moreover the image of the orthogonal complement \((\ell(u))^{\circ}\) under \(u\) is contained in \((\ell(u^*))^{\circ}\) . If \(u_1\) is the restriction of \(u\) to \((\ell(u))^{\circ}\) , considered as a mapping from \((\ell(u))^{\circ}\) into \((\ell(u^*))^{\circ}\) , the adjoint \(u_1^*\) is the restriction of \(u^*\) to \((\ell(u^*))^{\circ}\) ; if \(\ell(u) \neq E\) , let \(\langle u\rangle\) , called the sub-norm of \(u\) , denote the norm \(\|u_1\|\) ; if \(\ell(u) = E\) , we put \(\langle u\rangle = 0\) . Then \(\langle u^*\rangle = \langle u\rangle\) .
\end{enumerate}

\begin{enumerate}
\def\labelenumi{\arabic{enumi})}
\setcounter{enumi}{5}
\tightlist
\item
  Let E be a hilbertian space, M, N two closed subspaces of E and \(M^{\circ}\) , \(N^{\circ}\) their respective orthogonal complements; let \(p_{M}\) , \(p_{N}\) denote the orthogonal projections on M and N respectively, such that \(1_{E} - p_{M}\) , \(1_{E} - p_{N}\) are the orthogonal projections on \(M^{\circ}\) and \(N^{\circ}\) respectively. Put \(u_{\mathrm{NM}} = (1_{\mathrm{E}} - p_{\mathrm{N}}) p_{\mathrm{M}}\) and \(\delta(\mathbf{M}, \mathbf{N}) = \|u_{\mathrm{NM}}\|\) ; we have \(\delta(\mathbf{M}, \mathbf{N}) = \delta(\mathbf{N}^{\circ}, \mathbf{M}^{\circ}) \leqslant 1\) ; the relation \(\delta(\mathbf{M}, \mathbf{N}) < 1\) implies \(M \cap N^{\circ} = \{0\}\) .
\end{enumerate}

\begin{enumerate}
\def\labelenumi{\alph{enumi})}
\item
  Let \(\tilde{\mathbf{M}}\) denote the orthogonal complement in \(\mathbf{M}\) of \(\mathbf{M} \cap \mathbf{N}^{\circ}\) , and let \(\varepsilon(\mathbf{M}, \mathbf{N}) = \delta(\tilde{\mathbf{M}}, \mathbf{N})\) . Show that (with the notations of exerc. 5) \(\ell(u_{\mathrm{NM}}) = \mathbf{M} \cap \mathbf{N}^{\circ}\) and deduce that \(\varepsilon(\mathbf{M}, \mathbf{N}) = \langle u_{\mathrm{NM}} \rangle \leqslant \delta(\mathbf{M}, \mathbf{N})\) ; moreover, if \(\mathbf{M} \cap \mathbf{N}^{\circ} = \{0\}\) (and in particular if \(\delta(\mathbf{M}, \mathbf{N}) < 1\) ), then \(\varepsilon(\mathbf{M}, \mathbf{N}) = \delta(\mathbf{M}, \mathbf{N})\) .
\item
  For a continuous linear mapping \(u\) from E into itself, we designate by conorm of \(u\) the number \(c(u) = \inf \|u(x)\|/\|x\|\) , where \(x\) ranges over the set of all vectors \(\neq 0\) orthogonal to \(u^{-1}(0)\) (if \(u = 0\) , put \(c(u) = 1\) ). For \(u(\mathrm{E})\) to be closed in E, it is necessary and sufficient that \(c(u) > 0\) (I, p.~17, th. 1). We have \(c(u^{*}) = c(u)\) .
\item
  Let \(v_{\mathbf{NM}} = p_{\mathbf{N}}p_{\mathbf{M}}\) . Show that
\end{enumerate}

\[
\varepsilon (\mathbf {M}, \mathbf {N}) ^ {2} + c (v _ {\mathbf {N M}}) ^ {2} = 1
\]

(observe that \(\|u_{\mathrm{NM}}(x)\|^{2} + \|v_{\mathrm{NM}}(x)\|^{2} = \|p_{\mathrm{M}} \cdot x\|^{2}\) , and that the kernel of \(v_{NM}\) is \(M^{\circ} + (M \cap N^{\circ})\) and deduce that \(\langle u_{NM}\rangle^{2} \leqslant 1 - c(v_{\mathrm{NM}})^{2}\) .)
\item
  Deduce from b) and c) that \(\varepsilon(\mathrm{N}, \mathrm{M}) = \varepsilon(\mathrm{M}, \mathrm{N})\) and, using a), that \(\varepsilon(\mathrm{M}^{\circ}, \mathrm{N}^{\circ}) = \varepsilon(\mathrm{M}, \mathrm{N})\) .
\item
  Put \(g(\mathbf{M}, \mathbf{N}) = \|p_{\mathbf{M}} - p_{\mathbf{N}}\|\) (cf.~V, p.~64, exerc. 17). Show that

\[
g (\mathbf {M}, \mathbf {N}) = \sup (\delta (\mathbf {M}, \mathbf {N}), \delta (\mathbf {N}, \mathbf {M}))
\]

(observe that \(p_{\mathbf{M}} - p_{\mathbf{N}} = (1_{\mathrm{E}} - p_{\mathbf{N}})p_{\mathbf{M}} - p_{\mathbf{N}}(1_{\mathrm{E}} - p_{\mathbf{M}})\) ); deduce that \(\varepsilon (\mathbf{M},\mathbf{N})\leqslant g(\mathbf{M},\mathbf{N})\) . If \(\mathbf{M}\cap \mathbf{N}^{\circ} = \mathbf{N}\cap \mathbf{M}^{\circ} = \{0\}\) , we have the relation \(\varepsilon (\mathbf{M},\mathbf{N}) = \delta (\mathbf{M},\mathbf{N}) = \delta (\mathbf{N},\mathbf{M}) = g(\mathbf{M},\mathbf{N})\) . If \(g(\mathbf{M},\mathbf{N}) < 1\) , then \(\mathbf{M}\cap \mathbf{N}^{\circ} = \mathbf{N}\cap \mathbf{M}^{\circ} = \{0\}\) .

\end{enumerate}

\begin{enumerate}
\def\labelenumi{\alph{enumi})}
\setcounter{enumi}{5}
\tightlist
\item
  Let \(Q_{M}\) , \(Q_{N}\) be two continuous projections in E, with respective images M and N; give another proof of the relation \(g(\mathrm{M}, \mathrm{N}) \leqslant \|Q_{\mathrm{M}} - Q_{\mathrm{N}}\|\) (V, p.~64, exerc. 17). (Note that for all \(x \in E\) , we have \(\|(1_{E} - Q_{M}).x\|^{2} + \|Q_{M}^{*}.x\|^{2} = \|x\|^{2} + \|(Q_{\mathrm{M}} - Q_{\mathrm{M}}^{*}).x\|^{2}\) , and apply this relation to \(x = (p_{\mathrm{M}} - p_{\mathrm{N}}).y\) , observing the relations \((1_{\mathrm{E}} - Q_{\mathrm{M}})(p_{\mathrm{M}} - p_{\mathrm{N}}) = (Q_{\mathrm{M}} - Q_{\mathrm{N}}) p_{\mathrm{N}}\) and \((p_{\mathrm{M}} - p_{\mathrm{N}}) Q_{\mathrm{M}} = (1_{\mathrm{E}} - p_{\mathrm{N}})(Q_{\mathrm{M}} - Q_{\mathrm{N}}).\) )
\end{enumerate}

¶ 7) a) With the notations of exerc. 6, show that, for \(\mathbf{M} + \mathbf{N}^{\circ}\) to be closed, it is necessary and sufficient that \(\mathbf{M} + \mathbf{N}^{\circ} = (\mathbf{M}^{\circ} \cap \mathbf{N})^{\circ}\) .

\begin{enumerate}
\def\labelenumi{\alph{enumi})}
\setcounter{enumi}{1}
\tightlist
\item
  Show that the following properties are equivalent :
\end{enumerate}

\(\alpha)\)

\[
\varepsilon (\mathbf {M}, \mathbf {N}) <   1;
\]

β) \(M + N^{\circ}\) is closed in E;

\(\gamma)\) If \(\tilde{\mathbf{M}}\) is the orthogonal complement of \(\mathbf{M} \cap \mathbf{N}^{\circ}\) in \(\mathbf{M}\) and \((\mathbf{N}^{\circ})^{\sim}\) that of \(\mathbf{M} \cap \mathbf{N}^{\circ}\) in \(\mathbf{N}^{\circ}\) , then \(\mathbf{E}\) is the direct sum of \(\mathbf{M}^{\circ} \cap \mathbf{N}\) , \(\mathbf{M} \cap \mathbf{N}^{\circ}\) , \(\tilde{\mathbf{M}}\) and \((\mathbf{N}^{\circ})^{\sim}\) . Moreover, if \(R\) and \(S\) are respectively the projections from \(\mathbf{E}\) onto \(\mathbf{M}\) and \((\mathbf{N}^{\circ})^{\sim}\) corresponding to this decomposition, then \(\| R \| = \| S \| = (1 - \varepsilon^2 (\mathbf{M}, \mathbf{N}))^{1/2}\) . (To see that \(\alpha\) ) implies \(\beta\) ) observe first that if \(x \in \tilde{\mathbf{M}}\) and \(y \in (\mathrm{N}^\circ)^{\sim}\) , then \(|\langle x | y \rangle| \leqslant \varepsilon (\mathbf{M}, \mathbf{N}) \| x \| . \| y \|\) ; then, let \(u = x + y + t\) be the decomposition of an element \(u \in \mathbf{M} + \mathrm{N}^\circ\) with \(x \in \mathbf{M}\) , \(y \in (\mathrm{N}^\circ)^{\sim}\) and \(t \in \mathbf{M} \cap \mathrm{N}^\circ\) , deduce that \(\| x \| \leqslant (1 - \varepsilon^2 (\mathbf{M}, \mathbf{N}))^{-1/2} \| u \|, \| y \| \leqslant (1 - \varepsilon^2 (\mathbf{M}, \mathbf{N}))^{-1/2} \| u \|\) and \(\| t \| \leqslant \| u \|\) . To prove that \(\gamma\) ) implies \(\alpha\) ), consider the decomposition \(v = v_1 + v_2\) for a \(v \in \tilde{\mathbf{M}}\) , where \(v_1\) is the orthogonal projection of \(v\) onto \(\tilde{\mathbf{N}}\) , the orthogonal complement of \(\mathbf{N} \cap \mathbf{M}^\circ\) in \(\mathbf{N}\) ; we have \(R.v_1 = v\) , and if we had \(\varepsilon (\mathbf{M}, \mathbf{N}) = 1\) , there would exist a sequence \((v_n) \in \mathbf{M}\) such that \(\| v_n \| = 1\) and such that \(\| (v_n)_1\|\) tends to 0. Next show that the restriction \(R_1\) of \(R\) to \(\tilde{\mathbf{N}}\) is a bijection from \(\tilde{\mathbf{N}}\) onto \(\tilde{\mathbf{M}}\) ; to calculate \(\| R \|\) , show that \(\| R_1^{-1} \| \leqslant (1 - \varepsilon^2 (\mathbf{M}, \mathbf{N}))^{1/2}\) .

\begin{enumerate}
\def\labelenumi{\alph{enumi})}
\setcounter{enumi}{2}
\tightlist
\item
  Deduce from b) that if \(M + N^{\circ}\) is closed, then so is \(M^{\circ} + N\) .
\end{enumerate}

\begin{enumerate}
\def\labelenumi{\arabic{enumi})}
\setcounter{enumi}{7}
\tightlist
\item
  \begin{enumerate}
  \def\labelenumii{\alph{enumii})}
  \tightlist
  \item
    Let E be a hilbertian space, and let T be a continuous linear mapping from E into itself such that \(\|T\|\leqslant1\) . Show that the relations \(T.x=x,\langle T.x|x\rangle=\|x\|\) , \(T^{*}.x=x\) are equivalent, and that the kernel of \(1_{E}-T\) and the closure of the image of \(1_{E}-T\) are the orthogonal complements.
  \end{enumerate}
\end{enumerate}

\begin{enumerate}
\def\labelenumi{\alph{enumi})}
\setcounter{enumi}{1}
\tightlist
\item
  Let \(T\) be a continuous linear mapping from \(E\) into itself, satisfying the inequality

\[
\| x - T. x \| ^ {2} \leqslant \| x \| ^ {2} - \| T. x \| ^ {2}\tag{1}
\]

for all \(x \in E\) . Then \(\|T.x\| < \|x\|\) for all x such that \(T.x \neq x\) ; for every \(x \in E\) , the sequence \((T^{n}.x)\) converges to a point P.x, and P is the orthogonal projector onto the kernel of \(1_{E} - T\) .
\item
  Let \(P_{1}, \ldots, P_{r}\) be orthoprojectors on E. Show that the product \(T = P_{1}P_{2} \ldots P_{r}\) satisfies relation (1) (argue by induction on r); hence the orthoprojector P is the orthogonal projector onto the intersection of the images of the projectors \(P_{j}\) (note that if \(\|P_{j}.x\| < \|x\|\) for an index j, then \(\|T.x\| < \|x\|\) ).
\end{enumerate}

\begin{enumerate}
\def\labelenumi{\arabic{enumi})}
\setcounter{enumi}{8}
\tightlist
\item
  Let E be a hilbertian space and \((P_{j})_{j\in\mathbb{N}}\) be a sequence of orthoprojectors an E, such that, for all \(j\in N\) , there exists an \(n_{j}\in N\) such that for all \(k\in N\) , at least one of the orthoprojectors \(P_{k}, P_{k+1}, ..., P_{k+n_{j}}\) is equal to \(P_{j}\) . Put \(R_{s}=P_{s}P_{s-1}\ldots P_{0}\) for all \(s\in N\) .
\end{enumerate}

\begin{enumerate}
\def\labelenumi{\alph{enumi})}
\tightlist
\item
  For every \(x\in E\) , let \(x_{s}=R_{s}\cdot x\) . Show that \(\sum_{s}\|x_{s-1}-x_{s}\|^{2}\leqslant\|x\|^{2}\) , and deduce that for
\end{enumerate}

every integer \(r \geqslant 1\) , \(x_{s + r} - x_s\) tends to 0 as \(s\) tends to \(+\infty\) .

\begin{enumerate}
\def\labelenumi{\alph{enumi})}
\setcounter{enumi}{1}
\item
  Let \((x_{s_k})\) be a sequence extracted from \((x_s)\) which tends to a limit \(y\) weakly; then every sequence \((x_{s_k + r})\) also tends to \(y\) weakly. Deduce that \(y\) belongs to each of the subspaces \(\mathbf{M}_j = P_j(\mathbf{E})\) . (For each \(j\) , there exists \(r_k\) such that \(0 \leqslant r_k \leqslant n_j\) and \(s_k + r_k = j\) ; show that the sequence \((x_{s_k + r_k})\) tends to \(y\) weakly.)
\item
  Show that the sequence \((x_{s})\) converges weakly to the orthogonal projection from \(x\) onto the intersection \(\mathbf{M}\) of the \(\mathbf{M}_{j}\) . (Reduce to the case where \(\mathbf{M} = \{0\}\) , and use \(b\) ) and the weak compactness of every closed ball in E.)
\end{enumerate}

¶ 10) Let E be a hilbertian space, u a positive endomorphism of E.

\begin{enumerate}
\def\labelenumi{\alph{enumi})}
\tightlist
\item
  Show that for all \(x \in E\) , we have

\[
\left\| u (x) \right\| ^ {2} \leqslant \| u \|. \langle u (x) | x \rangle
\]

(observe that \(\langle u(x)|u(x)\rangle^2\leqslant \langle u(x)|x\rangle \langle u^2 (x)|u(x)\rangle\) , by V, p.~3, prop. 2).
\item
  Let M be a closed vector subspace of E, \(\mathbf{M}^{\circ}\) its orthogonal complement. Let \(x\in \mathbf{M}\) , and let \(f(x)\) be the lower bound of \(\langle u(x + y)|x + y\rangle\) as \(y\) ranges over \(\mathbf{M}^{\circ}\) . For every \(\varepsilon >0\) , let \(\operatorname {E}(x,\varepsilon)\) be the set of all \(y\in \mathbf{M}^{\circ}\) such that \(\langle u(x + y)|x + y\rangle \leqslant f(x) + \varepsilon\) . Show that \(\operatorname {E}(x,\varepsilon)\) is convex and that for all \(z\in \mathbf{M}^{\circ}\) , we have

\[
\langle u (x + y) | z \rangle^ {2} \leqslant \varepsilon \langle u (z) | z \rangle \quad \text { for } \quad y \in \mathrm{E} (x, \varepsilon)\tag{*}
\]

(consider the function \(g:t \mapsto \langle u(x + y + tz)/x + y + tz \rangle\) of the real variable t, which attains its minimum at a point \(t_{0}\) and note that \(g(t_{0}) \geqslant f(x)\) and \(g(0) \leqslant f(x) + \varepsilon\) ).
\item
  For every integer \(n \geqslant 1\) , let \(y_{n} \in \mathrm{E}(x, 1/n)\) ; show that the sequence \((u(x + y_{n}))\) tends to a limit \(x_{1}\) belonging to M, and that the sequence \((\langle u(y_{n})|y_{n}\rangle)\) is bounded (find a bound for the numbers \(\langle u(y_{n} - y_{m})|y_{n} - y_{m}\rangle\) for \(m \geqslant n\) and \(|\langle u(x + y_{n})|z\rangle|\) for \(z \in M^{\circ}\) using inequality (*)).

\item
  Let \((y_n')\) be a sequence of points of \(\mathbf{M}^\circ\) such that \(\langle u(y_n' - y_m')|y_n' - y_m'\rangle\) is arbitrarily small when \(m\) and \(n\) are large enough, and such that the sequence \((u(x + y_n'))\) has a limit \(x_1' \in \mathbf{M}\) ;

show that \(x_1' = x_1\) . (Let \(Q(z) = \langle u(z)|z \rangle\) ; first show that the number \(Q((y_p - y_p') - (y_q - y_q'))\) is arbitrarily small as soon as \(p\) and \(q\) are large enough, and deduce that the sequence \((Q(y_n - y_n'))\) is bounded; using the fact that \(\langle x_1' - x_1|y_p - y_p' \rangle = 0\) for all \(p\) , show that the sequence \((Q(y_n - y_n'))\) tends to 0 and use \(a\) ).)

\item
  Deduce from \(d\) ) that the point \(x_{1}\) does not depend on the choice of the \(y_{n} \in \mathrm{E}(x, 1/n)\) , and that if we put \(u_{1}(x) = x_{1}\) , then \(u_{1}\) is a linear mapping from \(\mathbf{M}\) into itself. Show that \(0 \leqslant \langle u_{1}(x)|x\rangle \leqslant \langle u(x)|x\rangle\) for all \(x \in \mathbf{M}\) and consequently that \(u_{1}\) is continuous and is an endomorphism of \(\mathbf{M}\) which is \(\geqslant 0\) . (Observe that \(\langle u(x + y_n)|y_n\rangle\) tends to 0 and \(\langle u(x + y_n)|x + y_n\rangle\) tends to \(f(x)\) .)
\item
  Let \(p_{M}\) be the orthoprojector with image M, and let \(u_{0} = u_{1} \circ p_{M}\) . We have \(0 \leqslant u_{0} \leqslant u\) , and M is stable under \(u_{0}\) , and the restriction of \(u_{0}\) to \(M^{\circ}\) is null. Show that \(u_{0}\) is the largest element in the family of all endomorphisms \(v \geqslant 0\) such that \(v \leqslant u\) , that M is stable under v and is such that the restriction of v to \(M^{\circ}\) is null.
\end{enumerate}

\begin{enumerate}
\def\labelenumi{\arabic{enumi})}
\setcounter{enumi}{10}
\tightlist
\item
  Let \(\mathrm{E}\) and \(\mathrm{F}\) be two hilbertian spaces. Show that for every element \(u\) in \(\mathrm{E} \hat{\otimes}_2 \mathrm{F}\) , there exists an orthonormal sequence \((e_n)\) in \(\mathrm{E}\) , an orthonormal sequence \((f_n)\) in \(\mathrm{F}\) and a sequence \((\lambda_n)\) of numbers \(\geqslant 0\) such that \(\sum_{n} \lambda_n^2 < +\infty\) and that \(u = \sum_{n} \lambda_n e_n \otimes f_n\) ; then \(\| u \|_2^2 = \sum_{n} \lambda_n^2\)
\end{enumerate}

(cf.~V, p.~55, th. 2 and p.~53, th. 1).

¶ 12) Let E be a real hilbertian space and V be a closed convex cone in E, with vertex 0, let V° be the polar cone of V (in E, identified canonically with its dual).

\begin{enumerate}
\def\labelenumi{\alph{enumi})}
\item
  Show that every point \(x \in \mathbf{E}\) can be written uniquely in the form \(x = x_{+} - x_{-}\) where \(x_{+} \in \mathbf{V}\) and \(x_{-} \in \mathbf{V}^{\circ}\) , and \(\langle x_{+}|x_{-}\rangle = 0\) .
\item
  For every facet F of V (II, p.~87, exerc. 3), F is either the point 0 or is a convex cone with vertex 0; the set of all \(y \in V^{\circ}\) which are orthogonal to F is a closed facet \(F'\) of \(V^{\circ}\) (but this is not the « dual facet » of F in the sense of II, p.~87, exerc. 6, the latter being empty).
\item
  Take for E the set of all Hilbert-Schmidt endomorphisms of a real hilbertian space, and for V the set of positive elements in E. Show that we have \(V^{\circ} = V\) and in this case interpret the result of a) (to see that \(V \subset V^{\circ}\) , use cor. 1 of V, p.~56).
\item
  Under the hypothesis of c), let \(v \in V\) ; the set L of all \(x \in H\) such that \(\langle v(x)|x \rangle = 0\) , or, which is the same, such that \(v(x) = 0\) (V, p.~77, exerc. 10) is a closed vector subspace of H, and the facet F of v in V is the closed set of all \(u \in V\) such that \(u(x) = 0\) for all \(x \in L\) ; it can be identified with the cone of all positive Hilbert-Schmidt endomorphisms of the hilbertian space \(L^{\circ}\) . Deduce that the projection from E onto the convex set F (V, p.~11) is identical with the orthogonal projector from E onto the closed vector subspace of E generated by F.
\end{enumerate}

¶ 13) Let E be a hilbertian space, G a subgroup of the group of all automorphisms of the hilbertian space structure of E. Let \(E^{G}\) be the closed vector subspace of E consisting of all vectors invariant under G, and p the orthoprojector from E onto \(E^{G}\) .

\begin{enumerate}
\def\labelenumi{\alph{enumi})}
\item
  Show that the orthogonal complement of \(\mathbf{E}^{\mathrm{G}}\) in \(\mathbf{E}\) is the closed vector subspace generated by the vectors \(s.x - x\) , where \(s\) ranges over \(\mathbf{G}\) and \(x \in \mathbf{E}\) .
\item
  Let \(\mathbf{H}\) be a non-empty closed convex subset of \(\mathbf{E}\) which is stable under \(\mathbf{G}\) . Show that the projection of 0 onto \(\mathbf{H}\) belongs to \(\mathbf{E}^{\mathrm{G}}\) .
\item
  Suppose that H is the closed convex envelope of the orbit of a point x of E, and let a be the projection of 0 onto H. Show that \(a = p(x)\) and that \(H \cap E^{G}\) reduces to the point a (« Birkhoff-Alaoglu th. »). (Observe that x - a is contained in the orthogonal complement of \(E^{G}\) .)
\item
  Suppose G is generated by an automorphism \(u\) of E. Show that \(p(x) = \lim_{n\to \infty}\frac{1}{n + 1}\sum_{j = 0}^{n}u^{j}(x)\)
\end{enumerate}

for all \(x \in \mathrm{E}\) (if \(y_{n} = \frac{1}{n + 1} \sum_{j=0}^{n} u^{j}(x)\) , note that the sequence \((y_{n})\) has a weak limit point \(a\) , and that \(u(a) = a\) , then use \(c\) ).

\begin{enumerate}
\def\labelenumi{\alph{enumi})}
\setcounter{enumi}{4}
\tightlist
\item
  Suppose \(G\) is the image of a homomorphism \(t \mapsto u_t\) from \(\mathbf{R}\) onto the group of automorphisms of \(E\) , such that for all \(x \in E\) , \(t \mapsto u_t \cdot x\) is a continuous mapping from \(\mathbf{R}\) into \(E\) . Show that

\(p(x) = \lim_{T \to \infty} \frac{1}{T} \int_{0}^{T} u_t \cdot x \, dt\)

for all \(x \in E\) .

\end{enumerate}

\begin{enumerate}
\def\labelenumi{\alph{enumi})}
\setcounter{enumi}{5}
\tightlist
\item
  Suppose that there exists an element \(x \neq 0\) of \(\mathbf{E}\) and a number \(\alpha\) such that \(0 < \alpha < 1\) and \(\| s.x - x\| \leqslant \alpha \| x\|\) for all \(s \in \mathbf{G}\) . Show that \(\mathrm{E}^{\mathrm{G}} \neq \{0\}\) (use \(c\) ).
\end{enumerate}

\begin{enumerate}
\def\labelenumi{\arabic{enumi})}
\setcounter{enumi}{13}
\tightlist
\item
  Let E be a complex hilbertian space, and T a Hilbert-Schmidt endomorphism of E. a) Let R and L be positive Hilbert-Schmidt endomorphisms such that \(R^{2} = T^{*}T\) and \(L^{2} = TT^{*}\) (V, p.~57, cor. 3); let \(R = \text{abs}(T)\) , and call this the « absolute value » of T (cf.~V, p.~75, exerc. 4); we have \(L = \text{abs}(T^{*})\) . Show that \(\text{Ker}(T) = \text{Ker}(R)\) and \(\overline{L(\text{E})} = \overline{T(\text{E})}\) . There exists one, and only one isometry V from \(R(\text{E})\) onto \(T(\text{E})\) such that T = VR; if we extend V by continuity to \(\overline{R(\text{E})}\) , then to an operator \(U \in \mathcal{L}(\text{E})\) by taking U.x = 0 on the orthogonal complement of \(R(\text{E})\) , we also have T = UR (polar decomposition of T). Then \(R = U^{*}T = U^{*}UR = RU^{*}U\) and \(L = URU^{*}\) , \(T = LU^{*}\) . If T belongs to \(\mathcal{L}^{1}(\text{E})\) , then so does \(R = \text{abs}(T)\) , and T is the product of two Hilbert-Schmidt endomorphisms.
\end{enumerate}

\begin{enumerate}
\def\labelenumi{\alph{enumi})}
\setcounter{enumi}{1}
\item
  If \(T\) belongs to \(\mathcal{L}^1 (\mathrm{E})\) , show that \(\operatorname {Tr}(\operatorname {abs}(T)) = \sup (\sum_i |\langle a_i | T.b_i\rangle |)\) where, on the right hand side, \((a_{i})\) and \((b_{i})\) range over the set of orthonormal bases of E (use the polar decomposition of \(T\) ). Show that if we put \(\| T\| _1 = \operatorname {Tr}(\operatorname {abs}(T))\) , then \(\| T\| _1\) is a norm on the space \(\mathcal{L}^1 (\mathrm{E})\) , such that \(\| T\| _2\leqslant \| T\| _1\) .
\item
  Conversely, if \(T \in \mathcal{L}(\mathrm{E})\) is such that, for every pair \(((a_i), (b_i))\) of orthonormal bases of \(\mathbf{E}\) , the sum \(\sum_{i} |\langle a_i | T \cdot b_i \rangle|\) is finite, then \(T \in \mathcal{L}^1(\mathrm{E})\) (first observe that \(T\) is a Hilbert-Schmidt endomorphism, then use the polar decomposition of \(T\) ).
\item
  Let \((\bar{T}_{\nu})\) be a sequence of Hilbert-Schmidt endomorphisms (resp. of elements of \(\mathcal{L}^1 (\mathrm{E})\) ) such that for every pair of points \(x,y\) of E, the sequence \((\langle x|T_{\nu}.y\rangle)\) converges to \(\langle x|T.y\rangle\) , where \(T\) is a linear mapping from E into itself; in addition, assume that the sequence of norms \(\| T_{\nu}\| _2\) (resp. \(\| T_{\nu}\| _1\) ) is bounded. Show that \(T\) is a Hilbert-Schmidt endomorphism (resp. an element of \(\mathcal{L}^1 (\mathrm{E})\) ) (use \(b\) ) and \(c\) ).
\item
  Deduce from \(d\) ) that the space \(\mathcal{L}^1 (\mathrm{E})\) is a Banach space for the norm \(\| T\| _1\) .
\item
  For an endomorphism \(T \in \mathcal{L}(\mathrm{E})\) to belong to \(\mathcal{L}^1(\mathrm{E})\) it is necessary and sufficient that, for at least one orthonormal basis \((e_i)\) of \(\mathrm{E}\) , the sum \(\sum_{i} \| T.e_i\|\) is finite (with the notations of \(a\) ), note that \(|\langle e_i|R.e_i\rangle| \leqslant \| T.e_i\|\) ).

\item
  In the space \(\ell_{\mathbf{C}}^2\) , let \((e_n)\) be the canonical orthonormal basis, and let \(a = \sum_{n=0}^{\infty} \frac{1}{n+1} e_n\) ;

if \(\mathrm{F}\) is the 1-dimensional subspace \(\mathbf{C}.a\) , the orthoprojector \(p_{\mathrm{F}}\) has finite trace, but the series \(\sum_{n}\| p_{\mathrm{F}}\cdot e_n\|\) is not convergent.
\end{enumerate}

\begin{enumerate}
\def\labelenumi{\arabic{enumi})}
\setcounter{enumi}{14}
\tightlist
\item
  Let \(\mathbf{E}\) be a complex hilbertian space; let \(\mathcal{B} = \mathcal{L}(\mathbf{E})\) denote the algebra of continuous endomorphisms of \(\mathbf{E}\) , endowed with the usual norm \(\| T\| = \sup_{\| x\| \leqslant 1}\| T.x\|\) . For every pair of points \(x, y\) in E, let \(\omega_{x,y}\) denote the continuous linear form \(T \mapsto \langle x|T.y \rangle\) on \(\mathcal{B}\) , and let \(\mathcal{B}_{\circ}\) be the closed subspace of the strong dual \(\mathcal{B}'\) of the Banach space \(\mathcal{B}\) , generated by the \(\omega_{x,y}\) . a) Show that the linear mapping which associates to every \(T \in \mathcal{B}\) the linear form \(\omega \mapsto \langle \omega, T \rangle\) on \(\mathcal{B}_{\circ}\) , is an isometry from \(\mathcal{B}\) onto the strong dual of \(\mathcal{B}_{\circ}\) ; in other words, \(\mathcal{B}_{\circ}\) is a predual (IV, p.~56, exerc. 23) of \(\mathcal{B}\) .
\end{enumerate}

The topology \(\sigma (\mathcal{B},\mathcal{B}_{\circ})\) on \(\mathcal{B}\) is called the ultraweak topology.

\begin{enumerate}
\def\labelenumi{\alph{enumi})}
\setcounter{enumi}{1}
\item
  For every element \(T\) of \(\mathcal{L}^1(\mathrm{E})\) , we define a linear form \(\phi_T\) on \(\mathcal{B}\) by the formula \(\phi_T(S) = \operatorname{Tr}(ST)\) for every operator \(S \in \mathcal{B}\) . Show that \(\phi_T\) is continuous and that the mapping \(T \mapsto \phi_T\) is an isometry from the Banach space \(\mathcal{L}^1(\mathrm{E})\) (exerc. 14, e)) onto the Banach space \(\mathcal{B}_{\circ}\) (first consider the case when \(T\) has finite rank).
\item
  Let \(\mathcal{B}_{\circ \circ}\) be the vector subspace of \(\mathcal{B}_{\circ}\) generated by the \(\omega_{x,y}\) (in such a way that \(\mathcal{B}_{\circ} = \overline{\mathcal{B}}_{\circ \circ} \) ). Show that \(\mathcal{B}_{\circ \circ}\) is barrelled (note that a subset of \(\mathcal{B}\) which is bounded for \(\sigma (\mathcal{B},\mathcal{B}_{\circ \circ})\) is bounded for the norm topology).
\item
  Let \(F_{n}\) be the subspace of \(\mathcal{B}_{\circ \circ}\) which is the image of the set of all endomorphisms of E of rank \(\leqslant n\) under the isometry defined in b). Show that \(F_{n}\) is nowhere dense in \(\mathcal{B}_{\circ \circ}\) and deduce that \(\mathcal{B}_{\circ \circ}\) is not a Baire space.
\end{enumerate}

\backmatter
\backmatterentry{Historical notes}

(chapters I to V)

(N.B. --- The roman letters refer to the bibliography at the end of this note.)

The general theory of topological vector spaces was founded in the period around the years 1920 to 1930. But the ground work had been under preparation since long by the study of numerous problems of functional Analysis; we cannot retrace the history of the subject without indicating, at least briefly, how the study of these problems slowly (particularly since the beginning of the 20th century) led the mathematicians to an awareness of the relationship between the questions being considered and the possibility of formulating them in a much more general manner, and applying to them uniform methods of solutions.

It can be said that the analogies between Algebra and Analysis, and the idea of considering functional equations (i.e.~where the unknown is a function) as « limiting cases » of algebraic equations have their origins in the infinitesimal Calculus, which in some sense was invented to generalize « from the finite to the infinite ». But the direct algebraic ancestor of the infinitesimal Calculus is the Calculus of finite differences (cf.~FVR, Historical note of chapters I, II, III, p.~54-58) and not the solution of general linear systems; it was only after the middle of the 18th century that the first analogies between the latter and the problems of differential Calculus made their appearance in the study of the equations of vibrating strings. We shall not enter into the details of the history of this problem here; but the constant reappearance of two fundamental ideas stands out, both of which are apparently due to D. Bernoulli. The first consists in considering the oscillation of the string as a « limiting case » of the oscillation of a system of n point masses as n increases indefinitely; we know that, later, this problem, for n finite, was the first example in the search for the eigenvalues of a linear transformation (cf.~A, Historical Notes of chapters VI-VII); these numbers correspond in the limiting case, to the frequencies of the « eigen oscillations » of the string, which where observed experimentally long before, and whose theoretical existence had been established (notably by Taylor) at the beginning of the century. This formal analogy, although hardly ever mentioned later ((I, b), p.~390) never seems to have been lost of sight of during the 19th century; but as we shall see, it acquired its full importance only around the years 1890-1900.

The other idea of D. Bernoulli (perhaps inspired by experimental facts) is the « superposition principle » according to which the most general oscillation of the string should be « decomposable » by superposition of the « eigen oscillations »; mathematically speaking, this means that the general solution of the equation of vibrating strings should have a series development as \(\sum_{n} c_{n} \phi_{n}(x, t)\) , where the \(\phi_{n}(x, t)\)

represent the eigen oscillations. We know that this principle was the starting point of a long battle on the possibility of developing an « arbitrary » function in a trigonometric series, a battle which was settled by the works of Fourier and of Dirichlet only in the first third of the 19th century. But even before this result was obtained, there were other examples of series development in « orthogonal » functions * : spherical functions, Legendre polynomials, and also various systems of the form ( \(e^{i\lambda_n x}\) ), where the \(\lambda_n\) are no longer multiples by the same number; these had already been introduced in the 18th century in oscillation problems, as also by Fourier and Poisson in the course of their researches on the theory of heat. Around 1830, Sturm (I) and Liouville (II) systematized all the phenomena observed in these various particular cases into a general theory of oscillations for functions of one variable; they considered the differential equation

\[
\frac {d}{d x} \left(p (x) \frac {d y}{d x}\right) + \lambda \rho (x) y = 0 \quad (p (x) > 0, \rho (x) > 0)\tag{1}
\]

with the boundary conditions

\[
\begin{array}{l l} k _ {1} y ^ {\prime} (a) - h _ {1} y (a) = 0 \\ k _ {2} y ^ {\prime} (b) + h _ {2} y (b) = 0 \end{array} (h _ {1} k _ {1} \neq 0, h _ {2} k _ {2} \neq 0, a <   b)\tag{2}
\]

and proved the following fundamental results :

\begin{enumerate}
\def\labelenumi{\arabic{enumi})}
\item
  the problem has a non-zero solution only if \(\lambda\) takes one of the values of a sequence \((\lambda_{n})\) of numbers > 0, tending to \(+\infty\) ;
\item
  for each \(\lambda_{n}\) , the solutions are multiples of the same function \(v_{n}\) , which may be assumed « normalized » by the condition \(\int_{a}^{b} \rho v_{n}^{2} dx = 1\) , and for \(m \neq n\) we have \(\int_{a}^{b} \rho v_{m} v_{n} dx = 0\) ;
\item
  every twice differentiable function \(f\) on \([a, b]\) which satisfies the boundary conditions (2), can be developed in a uniformly convergent series as \(f(x) = \sum_{n} c_{n} v_{n}(x)\) ,
\end{enumerate}

where \(c_{n} = \int_{a}^{b}\rho fv_{n}dx\)

\begin{enumerate}
\def\labelenumi{\arabic{enumi})}
\setcounter{enumi}{3}
\tightlist
\item
  the equality \(\int_{a}^{b}\rho f^2 dx = \sum_n c_n^2\) holds (this equality had already been proved by Parseval in 1799, though in a purely formal manner, for the system of trigonometric functions; and from it « Bessel's inequality » follows immediately; the latter inequality was announced by Bessel (again for trigonometric series) in 1828).
\end{enumerate}

Half a century later, these properties were completed by the work of Gram (III) who, following the researches of Tchebichef, threw light on the relationship between the development in series of orthogonal functions and the problem of « best quadratic approximation » (a direct outcome of the « method of least squares » of Gauss in the theory of errors); the latter consists of the following : given a finite sequence of functions \((\psi_{i})_{1\leqslant i\leqslant n}\) , and a function f, to find the linear combination \(\sum_{i}a_{i}\psi_{i}\) for which the integral \(\int_{a}^{b}\rho (f - \sum_{i}a_{i}\psi_{i})^{2}dx\) attains its minimum. In principle, this only suggests a trivial linear algebraic problem, but Gram solved it in an original way, by applying the method of « orthonormalization » to the \(\psi_{i}\) , as described in chap.~V, p.~23 (and generally known under the name of Erhard Schmidt). Next, in the case of an infinite orthonormal system, the question arises of finding out when the « best quadratic approximation » \(\mu_{n}\) of a function f, by linear combinations of the first n functions of the sequence, tends to 0 as n increases indefinitely *; Gram was thus led to the definition of the notion of a complete orthonormal system, and recognized that this property is equivalent to the non-existence of non-zero functions which are orthogonal to all the \(\phi_{n}\) . He even attempted to elucidate the concept of « mean quadratic convergence », but before the introduction of the fundamental ideas of measure theory, he could hardly obtain any general results in this direction.
* This term however does not appear before the work of Hilbert.

In the second half of the 19th century, the major effort of analysts was mainly directed towards the extension of the Sturm-Liouville theory to functions of several variables. This theory was prompted by the study of elliptic partial differential equations arising in Mathematical Physics, and the boundary value problems which are naturally associated with these. The main interest primarily centered on the equation of « vibrating membranes »

\[
\mathrm{L} _ {\lambda} (u) \equiv \Delta u + \lambda u = 0\tag{3}
\]

where solutions vanishing on the boundary of a sufficiently regular domain G were sought; the methods which had worked successfully for functions of one variable were no longer appropriate for this problem, and the considerable analytic difficulties that presented themselves were overcome little by little. We recall the main steps towards the solution: the introduction of the « Green's function » of G, whose existence was proved by Schwarz; the proof, again due to Schwarz, of the existence of the smallest eigenvalue; and finally, in 1894, H. Poincaré, in a celebrated memoir (V a) succeeded in proving the existence and the essential properties of all the eigenvalues. He considered the solution of the equation \(\mathbf{L}_{\lambda}(u) = f\) , for a « second member » f given; the solution being such as to vanish on the boundary; then by a skillful generalization of Schwarz's method, he proved that \(u_{\lambda}\) is a meromorphic function of the complex variable \(\lambda\) , having only real simple poles \(\lambda_{n}\) , and these are precisely the eigenvalues being sought.

* It must be pointed out that in this study, Gram did not restrict himself to considering only continuous functions, but emphasised the importance of the condition \(\int_{a}^{b} \rho f^2 dx < +\infty\) .

These researches are directly related to the beginnings of the theory of linear integral equations, which must have certainly contributed the maximum to the advent of modern ideas. We shall here limit ourselves to giving a brief outline of the development of this theory (for fuller details, we refer to the Historical Notes which will follow the chapters of this Treatise dedicated to spectral theory). This kind of functional equations, which first made a modest appearance in the first half of the 19th century (Abel, Liouville), had already acquired some importance since Beer and C. Neumann reduced the solution of « Dirichlet's problem » for a sufficiently regular domain G to the solution of an « integral equation of second kind »

\[
u (x) + \int_ {a} ^ {b} \mathrm{K} (x, y) u (y) d y = f (x)\tag{4}
\]

for the unknown function u; C. Neumann succeeded in solving this equation in 1877 by a method of « successive approximations ». Prompted as much by the algebraic analogies mentioned above as by the results he had obtained for the equation of vibrating membranes, H. Poincaré, in 1896 (V b) introduced a variable parameter \(\lambda\) in front of the integral in the preceding equation, and asserted that, just as in the case of the equation of vibrating membranes, the solution is once again a meromorphic function of \(\lambda\) ; but he was unable to prove this result. This was established seven years later by I. Fredholm (VI) (for a continuous « kernel » K and a finite interval \([a, b]\) ). The last mentioned author, perhaps with a greater awareness than his predecessors, let himself be guided by the analogy of (4) with the linear system

\[
\sum_ {q = 1} ^ {n} \left(\delta_ {p q} + \frac {1}{n} a _ {p q}\right) x _ {q} = b _ {p} (1 \leqslant p \leqslant n)\tag{5}
\]

to obtain the solution of (4) as the quotient of two expressions, based on the model of determinants, which arise in Cramer's formulas. This, however was not a new idea: since the beginning of the 19th century, the method of « indeterminate coefficients » (which consists of obtaining an unknown function, assumed to have a series development \(\sum_{n} c_{n} \phi_{n}\) , where the \(\phi_{n}\) are known functions, by calculating the coefficients \(c_{n}\) ) had led to « linear systems with infinitely many unknowns »

\[
\sum_ {j = 1} ^ {\infty} a _ {i j} x _ {j} = b _ {i} (i = 1, 2, \dots).\tag{6}
\]

Fourier, who encountered such a system, still solved it like an 18th century mathematician: he suppressed all the terms with indices \(i\) or \(j\) greater than \(n\) , explicitly solved the finite system so obtained by Cramer's formulas, then passed to the limit by letting \(n\) tend to \(+\infty\) in the solution! Much later, when this jugglery was no longer acceptable, it was again by the theory of determinants that the problem was attacked; Since 1886 (following the work of Hill), H. Poincaré, then H. von Koch, had set up a theory of « infinite determinants», which permits the resolution of certain kinds of systems (6) by following the classical model; and in spite of the fact that these results were not directly applicable to the problem tackled by Fredholm, it is beyond doubt that von Koch's theory in particular, served as a model for the construction of Fredholm's « determinants ».

It was at this moment that Hilbert entered the scene and gave a new impetus to the theory (VII). To begin with, he completed the work of Fredholm by effectively carrying out the passage to the limit, which leads to the solution of (4) from that of (5); but he immediately brought in the corresponding passage to the limit in the theory of real quadratic forms, which arose automatically from integral equations with a symmetric kernel (i.e.~such that \(\mathbf{K}(y,x)=\mathbf{K}(x,y)\) ). These are the equations by far the most frequent in Mathematical Physics. He thus succeeded in obtaining the fundamental formula which directly generalizes the reduction of a quadratic form to its axes

\[
\int_ {a} ^ {b} \int_ {a} ^ {b} \mathrm{K} (s, t) x (s) x (t) d s d t = \sum_ {n = 1} ^ {\infty} \frac {1}{\lambda_ {n}} \left(\int_ {a} ^ {b} \phi_ {n} (s) x (s) d s\right) ^ {2},\tag{7}
\]

where the \(\lambda_{n}\) are the eigenvalues (necessarily real) of the kernel K, the \(\phi_{n}\) forming the orthonormal system of the corresponding eigen functions, and the second member of formula (7) is a convergent series if \(\int_{a}^{b} x^{2}(s) ds \leqslant 1\) . He also showed how every function which is « representable » as \(f(x) = \int_{a}^{b} \mathbf{K}(x, y) g(y) dy\) has a « development » \(\sum_{n=1}^{\infty} \phi_{n}(x) \int_{a}^{b} \phi_{n}(y) f(y) dy\) , and, following the analogy with the classical theory of quadratic forms, he indicated a procedure for determining the \(\lambda_{n}\) by a variational method. This is precisely the extension of the well-known extremal properties of the principal axes of a quadric surface ((VII), p.~1-38).

These preliminary results of Hilbert were almost immediately taken up by E. Schmidt, under a simpler and more general form, avoiding the introduction of « Fredholm's determinants » and also the passage from the finite to the infinite. The presentation was already very close to being abstract, the fundamental properties of linearity and of positivity of the integral being clearly the only facts used in the proof (VII \(a\) ). But by then Hilbert had developed much more general concepts. All the earlier works brought out the importance of square integrable functions, and Parseval's formula established a direct link between these functions and sequences ( \(c_n\) ) such that \(\sum_{n} c_n^2 < \infty\) . It is certainly this idea which guided Hilbert in his 1906 memoirs ((VII), chap.~XI-XIII), where, taking up the old method of « indeterminate coefficients » once again, he showed that the solution of the integral equation (4) is equivalent to the solution of an infinite system of linear equations

\[
x _ {p} + \sum_ {q = 1} ^ {\infty} k _ {p q} x _ {q} = b _ {p} (p = 1, 2, \dots)\tag{8}
\]

for the « Fourier coefficients » \(x_{p} = \int_{a}^{b} u(t) \omega_{p}(t) dt\) of the unknown function \(u\) with respect to a given complete orthonormal system \((\omega_{n})\) (with \(b_{p} = \int_{a}^{b} f(t) \omega_{p}(t) dt\) and \(k_{pq} = \int_{a}^{b} \int_{a}^{b} K(s, t) \omega_{p}(s) \omega_{q}(t) ds dt\) ). Moreover, from this point of view, the only solutions of (8) of interest are those for which \(\sum_{n} x_{n}^{2} < +\infty\) ; also it was to this kind of solution that Hilbert systematically restricted himself; but on the other hand, he extended the conditions imposed on the « infinite matrix » \(k_{pq}\) (which in (8) is such that \(\sum_{p,q} k_{pq}^{2} < +\infty\) ). Thenceforth, it was clear that the « Hilbert space » of all sequences \(x = (x_{n})\) of real numbers such that \(\sum_{n} x_{n}^{2} < +\infty\) , while not explicitly introduced was the space underlying the entire theory, and appears as a « passage to the limit » from a finite dimensional Euclidean space. In addition, and this was particularly important for later developments, Hilbert was led to introduce, not just one, but two distinct notions of convergence in this space (corresponding to what has since been called the weak topology and the strong topology *), as also a « principle of choice » which is precisely the property of weak compactness of the unit ball. The new linear algebra that he developed in connection with the solution of the system (8) depended entirely on these topological ideas: linear mappings, linear forms and bilinear forms (associated with linear mappings) were classified and studied with respect to their « continuity » properties **. In particular, Hilbert discovered that the success of Fredholm's method depended on the notion of « complete continuity », which he redeemed by formulating it for bilinear forms *** and by studying it profoundly; for more details we refer to the part of this Treatise where this important notion shall be developed, and also to the admirable and profound works of Hilbert, where he inaugurated the spectral theory of symmetric bilinear forms (bounded or not).

The language of Hilbert still remained classical, and throughout the « Grundzüge », he never lost sight of the applications of the theory which he developed from numerous examples (taking up almost half of the volume). The next generation already adopted a much more abstract point of view. Under the influence of the ideas of Fréchet and of F. Riesz on general topology (see Historical Notes of GT, chap.~I), E. Schmidt (VII b) and Fréchet himself, in 1907-1908 deliberately introduced the language of Euclidean geometry into the « Hilbert space » (real or complex); it is in these works that we find the first mention of the norm (with its present notation \(\| x \|\) ), the triangle inequality that it satisfies, and the fact that a Hilbert space is « separable » and complete; in addition, E. Schmidt proved the existence of the orthogonal projection onto a closed linear variety, which allowed him to give a simpler and more general form to Hilbert's theory of linear systems. Also in 1907, Fréchet and F. Riesz observed that the space of square integrable functions has an analogous « geometry », an analogy which was perfectly explained when, a few months later, F. Riesz and E. Fischer proved that this space is complete and is isomorphic to a « Hilbert space », and at the same time displayed in a striking manner the value of the tool newly created by Lebesgue. From this moment onwards, the essential points of the theory of hilbertian spaces could be considered to have been achieved. Among the later developments the axiomatic presentation of the theory by M. H. Stone and J. von Neumann around 1930 must be mentioned, and also the removal of the restrictions of « separability » which was the result of the work of Rellich, Löwig and F. Riesz (IX e) around the year 1934.
* The Calculus of variations had naturally led to different notions of convergence on the same set of functions (according to the requirement of uniform convergence of functions, or of uniform convergence of functions and of a certain number of their derivatives); but the modes of convergence defined by Hilbert were entirely new at that time.
** It must be pointed out that until around 1935, by a « continuous » function it was generally meant that this was a mapping wich transformed every convergent sequence into a convergent sequence.
*** For Hilbert, a bilinear form \(B(x, y)\) was completely continuous if, whenever the sequences \((x_{n})\), \((y_{n})\) tended weakly to \(x\) and \(y\) respectively, \(B(x_{n}, y_{n})\) tended to \(B(x, y)\).

Meanwhile, in the first few years of the 20th century, other streams of ideas came and reinforced the trend which led to the theory of normed spaces. The general idea of « functional » (i.e.~a numerical function defined on a set whose elements are themselves numerical functions of one or of several real variables) was redeemed in the last decades of the 19th century in connection with the calculus of variations on the one hand, and on the other, with the theory of integral equations. But it was primarily from the Italian school, around Pincherle, and above all Volterra, that the general idea of « operator » arose. The works of this school often stayed at a rather formal level and were related to particular problems, for lack of a sufficiently deep analysis of the underlying topological concepts. In 1903, Hadamard inaugurated the modern theory of « topological » duality, in his search for the most general continuous linear « functionals » on the space \(\mathcal{C}(\mathrm{I})\) of continuous numerical functions on a compact interval (endowed with the topology of uniform convergence), and he characterized these as limits of sequences of integrals \(x\mapsto \int_{\mathrm{I}}k_n(t)x(t)dt\) .

In 1907, Fréchet and F. Riesz proved similarly that the continuous linear forms on a Hilbert space are the « bounded » linear forms introduced by Hilbert; then in 1909, F. Riesz put Hadamard's theorem in a definitive form by expressing every continuous linear form on \(\mathcal{C}(\mathrm{I})\) as a Stieltjes integral, a theorem which much later served as the starting point for the modern theory of integration (see Historical Notes of INT, chap.~II-V).

The following year, F. Riesz (IX a) again made new and important progress in the theory by introducing and studying (modelled on the theory of the Hilbert space) the space \(L^{p}(I)\) of functions on an interval I whose p-th power is integrable (for an exponent p such that \(1 < p < +\infty\) ); three years later, this study was followed by analogous work on the sequence spaces \(\ell^{p}(\mathbf{N})\) (IX c). These researches, as we shall see later, made a great contribution towards the classification of ideas on duality, in the sense that for the first time we encountered two spaces in duality which were not naturally isomorphic *.

Then onwards, F. Riesz thought of an axiomatic study which would encompass all these results ((IX a), p.~452), and it seems that only the scruples of an analyst anxious not to deviate from classical mathematics restrained him from writing his celebrated memoire of 1918 on Fredholm's theory (IX d) in this form. There he mainly considered the space \(\mathcal{C}(\mathrm{I})\) of continuous functions on a compact interval; but after defining the norm of this space, and having remarked that \(\mathcal{C}(\mathrm{I})\) endowed with this norm is complete, he did not use anything other than the axioms of complete normed spaces in his arguments **. Without entering into a detailed examination of this work, we mention that the notion of a completely continuous mapping was defined (by the property of transforming a neighbourhood into a relatively compact set) in a general way for the first time in this work ***; by a masterpiece of axiomatic analysis, the entire theory of Fredholm (with respect to its qualitative aspect) was reduced to a single fundamental theorem, that every locally compact normed space is finite dimensional.
* In spite of the fact that the duality between \(L^1\) and \(L^\infty\) was implicit in most of the works of this epoch on the Lebesgue integral, it was only in 1918 that H. Steinhaus proved that every continuous linear form on \(L^1(\mathrm{I})\) (I a finite interval) is of the form \(x\mapsto\int_{\mathrm{I}} f(t)x(t)dt\), where \(f\in L^\infty(\mathrm{I})\).
** F. Riesz however, explicitly noted that the applications of his theorems to continuous functions is only a « touchstone » of much more general concepts (IX d), p.~71).
*** In his work on \(L^p\) spaces, F. Riesz had defined completely continuous mappings as those which transform every weakly convergent sequence into a strongly convergent sequence; this (on account of the weak compactness of the unit ball in the \(L^p\) for \(1 < p < +\infty\)) is equivalent to the above definition in this case; in addition, F. Riesz indicated that for the \(L^2\) spaces, his definition was equivalent to that of Hilbert (by translating from the language of linear mappings to that of bilinear forms (IX a), p.~487)).

The general definition of normed spaces was given in 1920-1922 by S. Banach, H. Hahn and E. Helly (the latter considered only sequence spaces of real or complex numbers). In the ten years that followed, the theory of these spaces developed mainly around two questions of fundamental importance for applications: the theory of duality and the theorems linked with the notion of Baire « category ».

We have seen that the idea of duality (in the topological sense) originated in the beginning of the 20th century; it was the underlying notion in Hilbert's theory and occupied a central place in the work of F. Riesz. The latter, for example, observed in 1911 ((IX b), p.~41-42), that the relation \(|f(x)| \leqslant M \| x\|\) (taken as the definition of « bounded » linear functionals in a Hilbert space) is equivalent to the continuity of \(f\) in the case of the space \(\mathcal{C}(\mathrm{I})\) , and this was proved by fairly general arguments. Concerning the characterization of continuous linear functionals on \(\mathcal{C}(\mathrm{I})\) , he further observed that the condition for a set A to be total in \(\mathcal{C}(\mathrm{I})\) is that there exist no Stieltjes measure \(\mu \neq 0\) on I which is « orthogonal » to all the functions in A (thus generalizing Gram's condition for complete orthonormal systems); finally, in the same work, he established that the dual of the space \(L^{\infty}\) is « bigger » than the space of Stieltjes measures ((IX b), p.~62).

On the other hand, F. Riesz, in his work on the spaces \(L^{p}(\mathbf{I})\) and \(\ell^{p}(\mathbf{N})\) succeeded in modifying the method of the solution of linear systems in a Hilbert space, as given by E. Schmidt (VIII b) so as to be applicable in more general cases. E. Schmidt's idea consisted in determining an « extremal » solution of (6) by finding a point in the closed linear variety represented by the equations (6), whose distance from the origin is the minimum. Using the same idea, F. Riesz showed that a necessary and sufficient condition such that there exists a function \(x \in \mathrm{L}^{p}(a, b)\) satisfying the equations

\[
\int_ {a} ^ {b} \alpha_ {i} (t) x (t) d t = b _ {i} (i = 1, 2, \dots)\tag{9}
\]

(where the \(\alpha_{i}\) belong to \(\mathbf{L}^q\) (with \(\frac{1}{p} +\frac{1}{q} = 1\) ), and such that in addition \(\int_{a}^{b}|x(t)|^{p}dt\leqslant \mathbf{M}^{p}\) , is that, for every finite sequence \((\lambda_i)_{1\leqslant i\leqslant n}\) of real numbers, we have

\[
\left| \sum_ {i = 1} ^ {n} \lambda_ {i} b _ {i} \right| \leqslant \mathrm{M} \left(\int_ {a} ^ {b} \left| \sum_ {i = 1} ^ {n} \lambda_ {i} \alpha_ {i} (t) \right| ^ {q} d t\right) ^ {1 / q}.\tag{10}
\]

In 1911 (IX b), he treated, in an analogous manner, the « problem of generalized moments », which consists of the solution of the system

\[
\int_ {a} ^ {b} \alpha_ {i} (t) d \xi (t) = b _ {i} (i = 1, 2, \dots)\tag{11}
\]

where the \(\alpha_{i}\) are continuous and the unknown is a Stieltjes measure \(\xi\) * ; it was clear in this case that the problem can be restated by saying that it consists in determining a continuous linear functional on \(\mathcal{C}(\mathrm{I})\) from its values on a given sequence of points in this space. It was in this form that Helly treated the problem in 1912 --- obtaining F. Riesz's conditions by a rather different method of much wider scope * --- and which he again took up in 1921, with much more general conditions. Introducing the notion of a norm (on the sequence spaces), as we have seen above, he observed that this notion generalizes that of the « gauge » of a convex body in an n-dimensional space, as used by Minkowski in his celebrated work on the « geometry of numbers » (IV). In the course of his researches, Minkowski also defined (in \(R^{n}\) ) the notions of a support hyperplane and of a « support function » (IV b), and proved the existence of a support hyperplane at every point of the boundary of a convex body ((IV a), p.~33-35). Helly extended these notions to a space of sequences E, endowed with an arbitrary norm; he established a duality between E and the space \(E'\) of sequences \(u = (u_{n})\) such that for all \(x = (x_{n}) \in E\) , the series \((u_{n}x_{n})\) is convergent; letting \(\langle u, x \rangle\) denote the sum of this series, he defined a norm in \(E'\) by the formula \(\sup_{x \neq 0} |\langle u, x \rangle| / \|x\|\) , which gives the support function in finite dimensional spaces **. Then Helly proved that the solution of a system (6) in E, where each sequence \(u_{i} = (a_{ij})_{j \geqslant 1}\) is assumed to belong to \(E'\) , reduces to the successive resolution of the following two problems: 1. to find a continuous linear form L on the normed space \(E'\) , such that \(\mathrm{L}(u_{i}) = b_{i}\) for every index i; this, as he pointed out, leads to conditions of the type (10); 2. to find if such a linear form can be written as \(u \mapsto \langle u, x \rangle\) for some \(x \in E\) . The latter problem, he observed, does not necessarily have a solution even if L exists, and he gave some sufficient conditions which imply the existence of the solution \(x \in E\) in some particular cases (X).
* The classical « problem of moments » corresponds to the case where the interval \(]a, b[\) is \(]0, +\infty[\) or \(]-\infty, +\infty[\), and where \(\alpha_{i}(t) = t^{i}\); moreover, one assumes that the measure \(\xi\) is positive (in his 1911 memoire, F. Riesz indicated how his general conditions must be modified when solutions of this nature are sought). Among the various methods for the solution of the classical problem of moments, we particularly mention that of F. Riesz, who very elegantly combined the general ideas of functional Calculus and the theory of functions of one complex variable to obtain explicit conditions on the \(b_{i}\). (Sur le problème des moments, 3, Ark. för Math., t. XVII (1922-1923), no 16, 52 p.)
* Like F. Riesz ((IX b), p.~49-50), Helly used a « principle of choice » in his proof, which is precisely the weak compactness of the unit ball in the space of Stieltjes measures; F. Riesz had also used the analogous property in the \(L^p\) spaces (1 < p < +\infty).
** To obtain a norm in this way, we must assume that the relation \(\langle u,x\rangle = 0\) for all \(x \in E\) implies \(u=0\), as is explicitly remarked by Helly.
*** Banach had already given an analogous argument in 1923 for defining an invariant measure in the plane (defined for every bounded subset) (XII a).

In 1927, these ideas were given their definitive form in a fundamental memoire of H. Hahn (XI), whose results were rediscovered (independently) by S. Banach two years later (XII b). Hahn applied the methods of Minkowski-Helly to an arbitrary normed space, and thus defined the structure of a (complete) normed space on the dual space; this immediately allowed Hahn to consider successive duals of a normed space, and to pose the problem of reflexive spaces in a general way, as already foreseen by Helly. But above all, the principal problem of the extension of a continuous linear functional without increasing its norm was definitely solved by Hahn in general, by an argument of transfinite induction on the dimension --- thus giving one of the first examples of an important application of the axiom of choice to functional Analysis ***. To these results, Banach added a detailed study of the relations between a continuous linear mapping and its transpose, extending to general normed spaces results previously known in the case of \(\mathbf{L}^p\) spaces only (IX a), by means of a deep theorem on weakly closed subsets of the dual (cf.~IV, p.~25, cor. 2); these results can be expressed in a more striking way using the notion of the quotient space of a normed space, which was introduced a few years later by Hausdorff and by Banach himself. Finally, it was once again Banach who discovered the relation between the weak compactness of the unit ball (observed in several particular cases, as mentioned above) and reflexivity, at least for the spaces satisfying the first axiom of countability. Since then, the broad outlines of the theory of duality of normed spaces could be considered to have been fixed.

During the same epoch some seemingly paradoxical theorems, whose first examples originated in the years around 1910, were clarified. In that year, Hellinger and Toeplitz had essentially proved that a sequence of bounded bilinear forms \(\mathbf{B}_{n}(x,y)\) on a Hilbert space, whose values \(\mathbf{B}_{n}(a,b)\) for every pair \((a,b)\) are bounded (by a number depending a priori on a and b) is in fact uniformly bounded on every ball. Their proof was based on an argument of reductio ad absurdum, by inductively constructing a particular pair \((a,b)\) violating the hypothesis; this method is since then known under the name « gliding bump », and is still useful in many analogous questions (cf.~IV, p.~54, exerc. 15). In 1905, Lebesgue had used a similar method to prove the existence of continuous functions whose Fourier series diverges at some points; and in the same year as Hellinger and Toeplitz he used the method again, to prove that a weakly convergent sequence in \(L^{1}\) is bounded in norm *. These examples multiplied in the following years, but without the introduction of any new ideas until 1927, when Banach and Steinhaus (with the partial collaboration of Saks) related these phenomena to the notion of a thin set and to Baire's theorem in complete metric spaces, thus obtaining a general assertion which encompassed all the previous particular cases (XIII). During the same epoch, the study of questions of « category » in complete normed spaces led Banach to several other results on continuous linear mappings; the most remarkable and certainly the deepest being the « closed graph » theorem which, like the Banach-Steinhaus theorem, has become a vital tool in modern functional Analysis (XII b).

The publication of Banach's treatise on « Linear Operators » (XII c) marks the coming of age for the theory of normed spaces. All the above mentioned results and many others can be found in this volume, though in a somewhat disorganized manner, but with many striking examples drawn from various domains of Analysis, and which seemed to forecast a brilliant future for the theory. The work had considerable success, and one of the immediate effects was the almost universal adoption of the language and notations used by Banach. But in spite of the great number of researches undertaken during the past 40 years on Banach spaces (XVII), if the theory of Banach algebras and its applications to commutative and non-commutative harmonic analysis are excluded, then the almost complete absence of new applications of the theory to the great problems of classical Analysis somewhat undermines the hopes based on it.

It was more in the sense of widening and of a deeper axiomatic analysis related to the concepts of normed spaces that the most fruitful developments took place.
* Observe also the analogous (easier) theorem proved by Landau in 1907 and which served as a starting point for F. Riesz in his theory of the \(L^p\) spaces: if the series with the general term \(u_n x_n\) converges for every sequence \((x_n)\in \ell^p(\mathbf N)\), then the sequence \((u_n)\) belongs to \(\ell^q(\mathbf N)\) (with \(\frac{1}{p}+\frac{1}{q}=1\)).







In spite of the fact that the functional spaces encountered since the beginning of the 20th century generally appear to be endowed with a « natural » norm, there are certain exceptions. Around 1910, E. H. Moore proposed a generalization of the notion of uniform convergence by replacing it with a notion of « relative uniform convergence », where a neighbourhood of 0 consists of functions \(f\) satisfying a relation \(|f(t)| \leqslant \varepsilon g(t)\) , \(g\) being a function which is everywhere \(>0\) and which could vary with the neighbourhood. On the other hand, before 1930, it was observed that notions such as simple convergence, convergence in measure for measurable functions, or compact convergence for entire functions, could not be defined by means of a norm; and in 1926, Fréchet observed that vector spaces of this kind could be metrizable and complete. But the theory of these more general spaces could be fruitfully developed only in relation with the idea of convexity. The latter (which already appeared in Helly's work) was the subject of the studies carried out by Banach and his students, who recognized the possibility of interpreting several results of the theory of normed spaces geometrically, thus preparing the road for a general definition of locally convex spaces, given by Kolmogoroff and J. von Neumann in 1935. The theory of these spaces and notably the questions related to duality, were mostly developed in the years 1950, and in this Book we have presented the essential results of this study. In this connection we must point out, on the one hand, the progress in simplicity and in generality, made possible by the focus on the fundamental concepts of general Topology developed between 1930 and 1940; and secondly, the importance of the notion of a bounded set, introduced by Kolmogoroff and von Neumann in 1935, and whose fundamental role in the theory of duality was brought to light by the work of Mackey (XIV) and of Grothendieck (XVIII). Finally, it is certain that the main impetus which motivated these researches came from the new possibilities of applications to Analysis in domains where Banach; theory did not work : in this connection, we mention the theory of sequence spaces developed since 1934, by Köthe, Toeplitz and their students in a series of memoirs (XV), the focus on the theory of « analytic functionals » of Fantappié, and above all, the theory of distributions by L. Schwartz (XVI), where the modern theory of locally convex spaces found a field of applications, which is certainly far from being exhausted.

\backmatterentry{Bibliography}

\begin{enumerate}
\def\labelenumi{(\Roman{enumi})}
\item
  C. STURM : a) Sur les équations différentielles linéaires du second ordre, Journ. de Math. (1), t. I (1836), p.~106-186 ; b) Sur une classe d'opérations à différences partielles, ibid., p.~373-444.
\item
  J. LIOUVILLE : a) Sur le développement des fonctions ou parties de fonctions en séries dont les divers termes sont assujettis à satisfaire à une même équation différentielle du second ordre contenant un paramètre variable, Journ. de Math. (1), t. I (1836), p.~253-265, t. II (1837), p.~16-35 et 418-436 ; b) D'un théorème dû à M. Sturm et relatif à une classe de fonctions transcendants, ibid., t. I (1836), p.~269-277.
\item
  J. P. GRAM, Ueber die Entwicklung reeller Functionen in Reihen mittelst der Methode der kleinsten Quadrate, J. de Crelle, t. XCIV (1883), p.~41-73.
\item
  H. MINKOWSKI : a) Geometrie der Zahlen, 1 \(^{re}\) éd., Leipzig (Teubner), 1896 ; b) Theorie der konvexen Körper, Gesammelte Abhandlungen, t. II, p.~131-229, Leipzig-Berlin (Teubner), 1911. (Réimpression, New York (Chelsea), 1967.)
\item
  H. POINCARÉ : a) Sur les équations de la Physique mathématique, Rend. Palermo, t. VIII (1894), p.~57-156 (= Œuvres, t. IX, p.~123-196, Paris (Gauthier-Villars), 1954); b) La méthode de Neumann et le problème de Dirichlet, Acta Mathematica, t. XX (1896), p.~59-142 (= Œuvres, t. IX, p.~202-272, Paris (Gauthier-Villars), 1954).
\item
  I. FREDHOLM, Sur une classe d'équations fonctionnelles, Acta Mathematica, t. XXVII (1903), p.~365-390.
\item
  D. HILBERT, Grundzüge einer allgemeinen Theorie der linearen Integralgleichungen, New York (Chelsea), 1953 (= Gött. Nachr., 1904, 1905, 1906, 1910).
\item
  E. SCHMIDT : a) Zur Theorie der linearen und nichtlinearen Integralgleichungen. I. Teil : Entwicklung willkürlicher Funktionen nach Systemen vorgeschriebener, Math. Ann., t. LXIII (1907), p.~433-476 ; b) Ueber die Auflösung linearer Gleichungen mit unendlich vielen Unbekannten, Rend. Palermo, t. XXV (1908), p.~53-77.
\item
  F. RIESZ : a) Untersuchungen über Systeme integrierbarer Funktionen, Math. Ann., t. LXIX (1910), p.~449-497 ; b) Sur certains systèmes singuliers d'équations intégrales, Ann. Ec. Norm. Sup. (3), t. XXVIII (1911), p.~33-62 ; c) Les systèmes d'équations linéaires à une infinité d'inconnues, Paris (Gauthier-Villars), 1913 ; d) Ueber lineare Funktionalgleichungen, Acta Mathematica, t. XLI (1918), p.~71-98 ; e) Zur Theorie des Hilbertschen Raumes, Acta litt. ac scient. (Szeged), t. VII (1934-35), p.~34-38.
\item
  E. HELLY, Ueber Systeme linearer Gleichungen mit unendlich vielen Unbekannten, Monatshefte für Math. und Phys., t. XXXI (1921), p.~60-91.
\item
  H. HAHN, Ueber lineare Gleichungssysteme in linearen Räumen, J. de Crelle, t. CLVII (1927), p.~214-229.
\item
  S. BANACH et H. STEINHAUS, Sur le principe de condensation des singularités, Fund. Math., t. IX (1927), p.~50-61.
\item
  S. BANACH : a) Sur le problème de la mesure, Fund. Math., t. IV (1923), p.~7-33 ; b) Sur les fonctionnelles linéaires, Studia Math., t. I (1929), p.~211-216 et 223-239 ; c) Théorie des opérations linéaires, Warszawa, 1932. (Réimpression, New York (Chelsea), 1963.)
\item
  G. W. MACKEY : a) On infinite-dimensional linear spaces, Trans. Amer. Math. Soc., t. LVII (1945), p.~155-207 ; b) On convex topological spaces, Trans. Amer. Math. Soc., t. LX (1946), p.~519-537.
\item
  G. KÖTHE, Neubegründung der Theorie der vollkommenen Räume, Math. Nachr., t. IV (1951), p.~70-80.
\item
  L. SCHWARTZ, Théorie des distributions, 2 \(^{e}\) édition, Paris (Hermann), 1966.
\item
  J. LINDENSTRAUSS and L. TZAFRIRI, Classical Banach spaces, t. I, Berlin-Heidelberg-New York (Springer), 1977.
\item
  A. GROTHENDIECK : a) Produits tensoriels topologiques et espaces nucléaires, Mem. Amer. Math. Soc., n° 16 (1955); b) Espaces vectoriels topologiques, 3e éd., São Paulo (Publ. Soc. Mat. São Paulo), 1964.
\end{enumerate}

\backmatterentry{Index of notation}

The reference numbers indicate the chapter and page (and, occasionally, exercise)

\(|\xi|\) , \(\|x\|\) : I, p.~3.

\(\mathcal{B}(I; K)\) , \(\mathcal{B}_{K}(I)\) , \(\ell_{K}^{\infty}(I)\) , \(\ell_{K}^{1}(I)\) , \(\mathcal{B}(I)\) , \(\ell^{1}(I)\) : I, p.~4.

\(E_{A}\) (A a convex symmetric set in a real vector space E): II, p.~26.

\(\langle x, y \rangle\) : II, p.~42.

\(\sigma(F, G)\) : II, p.~42.

\(M^{\circ}\) , \(M^{\circ\circ}\) : II, p.~44.

\(^{t}u\) (u a linear mapping): II, p.~46.

\(\mathcal{R}(X)\) : III, p.~9.

\(\mathcal{C}^{\infty}(U)\) : III, p.~9.

\(\mathcal{C}_{H}^{\infty}(U)\) , \(\mathcal{C}_{0}^{\infty}(U)\) : III, p.~9.

\(\mathcal{G}_{s,M}(I)\) , \(\mathcal{G}_{s}(I)\) , \(\mathcal{C}(I)\) : III, p.~10.

\(\mathcal{H}(U)\) , \(\mathcal{H}(L)\) (U an open subset of \(C^{n}\) , L a compact subset of \(C^{n}\) ): III, p.~10.

\(\mathcal{L}(E; F)\) : III, p.~13.

\(\mathcal{L}_{\Xi}(E; F)\) : III, p.~13.

\(\mathcal{L}_{s}(E; F)\) , \(\mathcal{L}_{c}(E; F)\) , \(\mathcal{L}_{pc}(E; F)\) , \(\mathcal{L}_{cc}(E; F)\) , \(\mathcal{L}_{b}(E; F)\) : III, p.~14.

\(E'\) , \(E'_{\Xi}\) , \(E'_{s}\) , \(E'_{c}\) , \(E'_{pc}\) , \(E'_{cc}\) , \(E'_{b}\) : III, p.~14.

\(\mathcal{L}(E)\) , \(\mathcal{L}_{\Xi}(E)\) , \(\mathcal{L}_{s}(E)\) , \(\mathcal{L}_{c}(E)\) , \(\mathcal{L}_{pc}(E)\) , \(\mathcal{L}_{cc}(E)\) , \(\mathcal{L}_{b}(E)\) : III, p.~14.

\(p_{M}\) (p a semi-norm, M a bounded subset): III, p.~14.

\(\mathcal{C}_{0}(\mathbf{R})\) : III, p.~18.

\(\tau(E, F)\) : IV, p.~2.

\(\beta(E, F)\) : IV, p.~4.

\(c_{E}\) : IV, p.~14.

\(\ell_{0}^{\infty}(\mathbf{N})\) : IV, p.~17.

\(c_{0}(\mathbf{N})\) , \(\ell^{1}(\mathbf{N})\) : IV, p.~18.

S(E): IV, p.~26.

\(H_{p}\) : IV, p.~26.

\(E_{\sigma}\) : IV, p.~32.

\(\mathcal{C}_{s}(\mathbf{X})\) : IV, p.~33.

\(C^{b}(\mathbf{X})\) , \(C(\mathbf{X})\) : IV, p.~36.

\(\mathcal{B}(\mathbf{X};\mathbf{R})\) : IV, p.~40.

Ind(u) (u a Fredholm operator): IV, p.~66, exerc. 21.

\(\overline{\xi}\) : V, p.~1.

\(\ell^{2}\) , \(\ell^{2}(\mathbf{N})\) : V, p.~4.

\(E_{(\mathbf{C})}\) : V, p.~4.

\(\langle x|y\rangle\), \(\|x\| = \langle x|x\rangle^{1/2}\), \((x|y) = \langle y|x\rangle\): V, p.~5.

\(\overline{\mathrm{E}}\) (E a complex prehilbertian space): V, p.~6.

\(\mathcal{H}^{s} (\text {Sobolev space}): V, p.~6.\) \(H^{2}(\text {D}): V, p.~7.\) \(\mathcal{C}_{0}^{1}(U): V, p.~8.\) \(p_{H} (\text {H convex separated and complete set in a prehilbertian space}): V, p.~10.\) \(x^{*} (x a vector of a Hilbert space): V, p.~15 et p.~40.\) \(\bigoplus_{i\in I}\mathrm{E}_{i},\bigoplus_{i\in I}\mathrm{E}_{i}:V,p.18.\) \(E_{1}\oplus E_{2}\oplus\cdots\oplus E_{n} (E_{i}\text { Hilbert spaces}):V,p.18.\) \(\ell_{E}^{2}(I),\ell^{2}(I)(E\text {a Hilbert space}):V,p.18.\)

\(E_{1} \otimes_{2} E_{2}, \|z\|_{2} (z \in E_{1} \otimes_{2} E_{2}) : V, p.~26.\)

\(E_{1} \otimes_{2} E_{2} \otimes_{2} \cdots \otimes_{2} E_{n}, \bigotimes_{i=1}^{n} E_{i}, \|z\|_{2} \left( z \in \bigotimes_{i=1}^{n} E_{i} \right) : V, p. 27.\)

\(E_{1} \hat{\otimes}_{2} E_{2} \hat{\otimes}_{2} \ldots \hat{\otimes}_{2} E_{n}, \bigotimes_{1 \leqslant i \leqslant n} \hat{\otimes}_{2} E_{i}: V, p. 28\)

\(u_{1}\hat{\otimes}_{2}u_{2}\hat{\otimes}_{2}\ldots \hat{\otimes}_{2}u_{n}(u_{i}\) linear mappings): V, p.~28.

\(\hat{\mathbf{T}}^n (\mathrm{E}),\mathrm{E}\hat{\otimes}_n:\mathrm{V},\) p.~29.

\(\hat{\mathbf{S}}^{n}(\mathrm{E}), \hat{\mathbf{S}}(\mathrm{E}) : \mathrm{V}, \mathrm{p.} 30.\)

\(\hat{\mathbf{A}}^{n}(\mathrm{E}), \hat{\mathbf{A}}(\mathrm{E}) : \mathrm{V}, \mathrm{p. 33}\) .

\(\hat{\mathbf{T}}^n (u),\hat{\mathbf{S}}^n (u)\) \((u\) a linear mapping):V,p.31 and p.32.

\(\hat{\Lambda}^n (u)\) ( \(u\) a linear mapping): V, p.~34.

v.u, vu (u, v linear mappings) : V, p.~37.

\(u^{*}\) (u a linear mapping): V, p.~38.

\(\mathcal{H}(\mathrm{E})\) (E a Hilbert space): V, p.~44.

\(u \geqslant 0\) (u an endomorphism of a Hilbert space): V, p.~45.

\(\mathcal{L}_{+}(\mathrm{E}):\mathrm{V},\mathrm{p.}45.\)

\(u \geqslant v\) (u, \(v\) in \(\mathcal{L}(\mathrm{E})\) , E a Hilbert space): V, p.~45.

\(\tau (u)\) ( \(u\) an endomorphism of finite rank): V, p.~48.

\(\operatorname{Tr}(u) (u \geqslant 0 \text{ in } \mathcal{L}(\mathrm{E})): \mathrm{V}, \text{ p. } 49.\)

\(\mathcal{L}^1 (\mathrm{E})\) (E a Hilbert space) : V, p.~51.

\(\mathcal{L}^2 (\mathrm{E};\mathrm{F}),\mathcal{L}^2 (\mathrm{E})\) (E, F Hilbert spaces): V, p.~52.

\(\| u\| _2\) \((u\in \mathcal{L}(\mathrm{E};\mathrm{F}),\mathrm{E},\mathrm{F}\) Hilbert spaces): V, p.~52.

\(\mathrm{Tr}(\mathrm{Q} / \mathrm{H})\) (Q, H positive quadratic forms): V, p.~57.

\backmatterentry{Index of terminology}

Absorbent set, absorption of one set by another : I, p.~7. Adapted bornology : III, p.~3. Adjoint : V, p.~38. Affine transformation : IV, p.~39. Associated (Hausdorff vector space) with a topological vector space : I, p.~4.

Balanced convex closed envelope of a set : II, p.~13 and p.~62. Balanced core of a set : I, p.~7. Balanced envelope of a set : I, p.~7. Balanced set : I, p.~6. Banach-Dieudonné theorem : IV, p.~24. Banach-Saks-Kakutani theorem : V, p.~68, exerc. 33. Banach space : I, p.~5. Banach-Steinhaus theorem : III, p.~25. Banach's theorem : I, p.~17. Barrel : III, p.~24. Barrelled space : III, p.~24. Base of a bornology : III, p.~1. Basis (algebraic) of a Hilbert space : V, p.~22. Basis (Banach) : IV, p.~69, exerc. 14. Basis (complete Banach, contracting Banach) : IV, p.~70, exerc. 15. Basis (orthonormal) : V, p.~22. Basis (unconditional Banach) : IV, p.~71, exerc. 16. Bessel's inequality : V, p.~21. Bidual : IV, p.~14. Bipolar theorem : II, p.~44. Birkhoff-Alaoglu theorem : V, p.~78, exerc. 13. Bishop-Phelps theorem : II, p.~77, exerc. 4. Bornological locally convex space : III, p.~12. Bornology : III, p.~1. Bornology (adapted) : III, p.~3. Bornology (canonical) : III, p.~3. Bornology (convex) : III, p.~2. Bornology generated by a family of sets : III, p.~1. Bornology (product) : III, p.~2. Bornivorous set : III, p.~39, exerc. 11. Bounded set : III, p.~2 and p.~37, exerc. 1.

Canonical bornology : III, p.~3. Canonical mapping of \(\bigoplus_{i\in I}E_i'\) in \(\left(\prod_{i\in I}E_i\right)'\) : IV, p.~13. Canonical mapping of E in E'\,' : IV, p.~14. Canonical mapping of E onto E' (E a Hilbert space) : V, p.~15. Canonical topology on a finite dimensional vector space : I, p.~2. Cap : II, p.~57. Cauchy-Schwarz inequality : V, p.~3. Closed graph theorem : I, p.~19. Closed half-spaces defined by a closed hyperplane : II, p.~15. Cobord : IV, p.~72, exerc. 3. 1-cocycle (continuous) : IV, p.~72, exerc. 3.

Compact linear mapping : III, p.~6. Compatible topology and structure of an ordered vector space : II, p.~15. Compatible vector space structure and preorder : II, p.~12. Compatible vector space structure and topology : III, p.~1. Compatible with the duality (locally convex topology) : IV, p.~1. Complement (orthogonal) : V, p.~13. Complete topological vector space : I, p.~5. Completion of a Hausdorff prehilbertian space : V, p.~8. Completion of a Hausdorff topological vector space : I, p.~6. Complex linear form : II, p.~61. Complex linear variety : II, p.~61. Complex locally convex space : II, p.~62. Complexification (prehilbertian space) : V, p.~5. Complexified topological vector space : II, p.~62. Concave function : II, p.~17. Cone (asymptotic) : II, p.~67. Cone (convex) generated by a set : II, p.~11. Cone (pointed and non-pointed) : II, p.~10. Cone (polyhedral) : II, p.~91. Cone (proper pointed convex) : II, p.~11. Conjugate of a complex prehilbertian space : V, p.~6. Convex balanced envelope of a set : II, p.~10 and p.~62. Convex bornology : III, p.~2. Convex closed envelope of a set : II, p.~13. Convex function : II, p.~17. Convex set : II, p.~7 and p.~62. Convex (symmetric) envelope of a set : II, p.~16. Coordinates with respect to an orthonormal base : V, p.~22. Core (balanced) of a set : I, p.~7.

Density of order : V, p.~7. Dimension (hilbertian) : V, p.~24. Dimension of a convex set : II, p.~10. Dirichlet space : V, p.~8. Distal set : IV, p.~72, exerc. 1. Distinguished space : IV, p.~52, exerc. 4. Dual (algebraic) of a real topological vector space : II, p.~42. Dual of a locally convex space (real or complex) : III, p.~14. Dual of a real topological vector space : II, p.~42. Dual (weak, strong) : III, p.~14. Duality separating in F, separating duality : II, p.~41. Duality (vector spaces in) : II, p.~40. Dvoretzky-Rogers theorem : V, p.~63, exerc. 14.

Eberlein's theorem : IV, p.~35. D. Edwards' theorem : II, p.~94, exerc. 41. Endomorphism (hermitian) : V, p.~44. Endomorphism (normal) : V, p.~43. Endomorphism (positive) : V, p.~45. Envelope of a set (balanced) : I, p.~7. Envelope of a set (balanced convex) : II, p.~10. Envelope of a set (balanced convex closed) : II, p.~13 and p.~62. Envelope of a set (closed convex) : II, p.~13. Envelope of a set (convex) : II, p.~9. Envelope of a set (symmetric convex) : II, p.~10. Envelope of a set (symmetric convex closed) : II, p.~13. S-equihypocontinuous, T-equihypocontinuous, (S, T)-equihypocontinuous set : III, p.~47, exerc. 7.

Filtered set of semi-norms : II, p.~3.

Family (topologically independent) : I, p.~11.

Exhaustion of a Hausdorff locally convex space, exhaustible space : III, p.~49, exerc. 1.

Extremal generator of a convex cone : II, p.~57.

Extremal point of a convex set : II, p.~54.

Facet : II, p.~87, exerc. 3.

Facet (dual) : II, p.~87, exerc. 6.

Family (orthonormal) : V, p.~21.

Final locally convex topology : II, p.~27.

Final subspace : V, p.~41.

Fock spaces : V, p.~32 and p.~34.

Form (bilinear) putting two spaces in duality : II, p.~40.

Form (complex linear, real linear) : II, p.~61.

Form (hermitian) : V, p.~1.

Form (positive hermitian) : V, p.~2.

Form (positive linear) : II, p.~13.

Form (separating hermitian) associated with a hermitian form : V, p.~2.

Fréchet space : II, p.~24 and p.~63.

Function (concave, convex, strictly concave, strictly convex) : II, p.~16-17.

Function (positive definite) : V, p.~8.

Function (positively homogeneous), sub-linear function : II, p.~19-20.

Fundamental system of semi-norms : II, p.~3.

Gauge of a convex set : II, p.~20.

Gaussian space : V, p.~32.

Generated (bornology) by a family of sets : III, p.~1.

Generator (extremal) of a convex cone : II, p.~57.

Gevrey's spaces : III, p.~10.

Gram's determinant : V, p.~71, exerc. 7.

Grothendieck's theorem : III, p.~20.

Group on which a mean can be defined : IV, p.~72, exerc. 4.

Haar's theorem : II, p.~83, exerc. 8.

Hadamard's inequalities : V, p.~37.

Hahn-Banach theorem : II, p.~22, p.~36 and p.~63.

Half-spaces (closed, open) determined by a closed hyperplane : II, p.~15.

Hardy space : V, p.~7.

Hausdorff completion of a topological vector space : I, p.~6.

Hausdorff vector space associated with a topological vector space : I, p.~4.

Helly's theorem : II, p.~68, exerc. 21.

Hermitian endomorphism : V, p.~44.

Hilbert-Schmidt mapping : V, p.~52.

Hilbert space, hilbertian space : V, p.~6.

Hyperplane (support) of a set : II, p.~37.

S-hypocontinuous, T-hypocontinuous, (S, T)-hypocontinuous bilinear mapping : III, p.~30.

Index of a Fredholm operator : IV, p.~66, exerc. 21.

Induced prehilbertian structure on a vector subspace : V, p.~6.

Inductive limit of locally convex spaces or topologies : II, p.~29.

Inequalities (Hadamard's) : V, p.~37.

Inequality (Bessel's) : V, p.~21.

Inequality (Cauchy-Schwarz) : V, p.~3.

Infra-barrelled space : III, p.~44, exerc. 7.

Initial subspace : V, p.~41.

Initial topology : IV, p.~4.

Internal point of a convex set : II, p.~26.

James-Klee theorem : IV, p.~57, exerc. 25. R. C. James' space : IV, p.~71, exerc. 18.

Krein-Milman theorem : II, p.~55. Krein's theorem : IV, p.~37.

Legendre polynomial : V, p.~24. Locally convex complex space : II, p.~62. Locally convex real space : II, p.~23. Locally convex topology : II, p.~23 and p.~62.

Mackey theorem : IV, p.~2. Mackey topology : IV, p.~2. Mapping (Hilbert-Schmidt) : V, p.~52. Markoff-Kakutani theorem : IV, p.~39. Matrix (Hilbert's) : V, p.~75, exerc. 3. Matrix with respect to orthonormal bases : V, p.~22. Mean : IV, p.~40. Metrisable topological vector space : I, p.~16. Minimal type (space of) : II, p.~85, exerc. 13. Montel space : IV, p.~18.

Normal endomorphism : V, p.~43.

Open half-spaces defined by a closed hyperplane : II, p.~15. Operator (Fredholm) : IV, p.~66, exerc. 21. Operator (unitary) : V, p.~41. Ordered vector space : II, p.~12. Orthogonal of a subspace for spaces in duality : II, p.~44. Orthogonal projector, orthoprojector : V, p.~13. Orthogonal sets in a prehilbertian space : V, p.~13. Orthogonal sets in spaces in duality : II, p.~41. Orthogonal to a subset in a prehilbertian space : V, p.~13. Orthogonal vectors for a duality : II, p.~41. Orthogonal vectors in a prehilbertian space : V, p.~5. Orthonormal basis : V, p.~22. Orthonormal set, family : V, p.~21. Orthonormalisation : V, p.~24. Orthoprojector (initial, final) : V, p.~41.

Parabolic convex set : II, p.~67, exerc. 17. Parseval's relation : V, p.~22. Partially isometric mapping : V, p.~42. Pointed cone : II, p.~10. Points (exposed, of strict convexity) : II, p.~88, exerc. 6. Points on the same side, strictly on the same side, of a hyperplane : II, p.~9. Polar decomposition of a Hilbert-Schmidt endomorphism : V, p.~79, exerc. 14. Polar of a set : II, p.~44 and 64. Polarization formulas : V, p.~2. Polyhedron : II, p.~90, exerc. 24. Positive endomorphism : V, p.~45. Positive hermitian form : V, p.~2. Positively homogeneous function : II, p.~19. Predual : IV, p.~56, exerc. 23. Prehilbertian semi-norm : V, p.~4. Prehilbertian space : V, p.~4. Preordered vector space : II, p.~12. Principle of condensation of singularities : III, p.~42, exerc. 10.

Product bornology : III, p.~2. Projection on a convex set : V, p.~11. Projection (orthogonal) : V, p.~13. Proper pointed convex cone : II, p.~11. Pythagoras' theorem : V, p.~12.

Quasi-complete space : III, p.~8.

Real linear form, real linear variety : II, p.~61. Real locally convex space : II, p.~23. Reflexive space : IV, p.~16. Relatively bounded space : III, p.~43, exerc. 6. Representation (unitary) : IV, p.~44. Ryll-Nardzewski theorem : IV, p.~43.

Scalar product : V, p.~5. Scalar square : V, p.~5. Segment (closed, open, open at x and closed at y) : II, p.~7. Semi-automorphism of prehilbertian spaces : V, p.~6. Semi-barrelled space : IV, p.~21. Semi-complete space : III, p.~7. Semi-norm : II, p.~1. Semi-normed space : II, p.~2. Semi-norm (prehilbertian) : V, p.~4. Semi-reflexive space : IV, p.~15. Separated completion of a prehilbertian space : V, p.~8. Separated completion of a topological vector space : I, p.~6. Separated (sets) by a closed hyperplane : II, p.~37. Separately continuous bilinear mapping : III, p.~28. Separately equicontinuous : III, p.~47. Separating duality : II, p.~41. Side of a hyperplane (points on the same, strictly on the same) : II, p.~8. Simplex : II, p.~71, exerc. 41. Šmulian theorem : IV, p.~36. Sobolev space : V, p.~6. Sole of a cone : II, p.~60. Space (DF) : IV, p.~57, exerc. 2. Space (topological vector) : I, p.~1. Space (weak) : II, p.~42. Square (scalar) : V, p.~5. Starshaped set : II, p.~65, exerc. 1. Strict inductive limit of a sequence of locally convex spaces, topologies : II, p.~33. Strictly concave, strictly convex function : II, p.~16-17. Strictly separated (sets) by a closed hyperplane : II, p.~37. Strong dual : III, p.~14. Strongly bounded subset of E' : III, p.~14. Sub-linear function : II, p.~20. Subspace (final, initial) of a continuous linear mapping : V, p.~41. Subspace (prehilbertian) : V, p.~6. Sum (external hilbertian) of Hilbert spaces : V, p.~18. Sum (hilbertian) of vector subspaces : V, p.~18. Sum (topological direct) of locally convex spaces or topologies : II, p.~30. Support variety : II, p.~87, exerc. 3. Symmetric convex closed envelope of a set : II, p.~13.

Tchebycheffs theorem : II, p.~84, exerc. 8. Tensor product (hilbertian) : V, p.~28. Tensor product of prehilbertian spaces : V, p.~26-27.

Topologically independent set, family : I, p.~12 and p.~11.

\(\mathfrak{S}\) -topology : III, p.~13 and IV, p.~2.

Topology compatible with an ordered vector space structure : II, p.~15.

Topology compatible with a vector space structure : I, p.~1.

Topology defined by a semi-norm, by a set of semi-norms : II, p.~2-3.

Topology (initial) : IV, p.~4.

Topology (locally convex) : II, p.~23 and p.~62.

Topology (Mackey) : IV, p.~2.

Topology of simple, compact, precompact, compact convex, bounded convergence : III, p.~14.

Topology (weak) : II, p.~42.

Topology (weakened) : IV, p.~4.

Total set : I, p.~11.

Trace of a continuous positive endomorphism : V, p.~49.

Trace of a positive quadratic form with respect to another : V, p.~57.

Transformation (affine) : IV, p.~39.

Transpose of a continuous linear mapping : II, p.~46 and IV, p.~6.

Ultrabornological space : III, p.~45, exerc. 19.

Ultranorm : I, p.~26, exerc. 12.

Ultra-semi-norm : II, p.~2.

Unitary operator : V, p.~41.

Unitary representation : IV, p.~44.

Unpointed cone : II, p.~10.

Variety (complex linear, real linear) : II, p.~61.

Variety (support) : II, p.~87, exerc. 3.

Weak dual : III, p.~14.

Weakly bounded : III, p.~14.

Weak topology, weak space : II, p.~42.

Weakened topology : IV, p.~4.

\backmatterentry{Summary of some important properties of Banach spaces}

For the reader's convenience, the principal results of normed spaces and, more particularly, of Banach spaces are collected here. The field of scalars K is either R or C.

Linear mapping spaces ; dual

\begin{enumerate}
\def\labelenumi{\arabic{enumi})}
\tightlist
\item
  Let \(\mathbf{E}\) and \(\mathbf{F}\) be two normed spaces. A linear mapping \(u\) of \(\mathbf{E}\) in \(\mathbf{F}\) is continuous, if and only if
\end{enumerate}

\[
\| u \| = \sup _ {\| x \| \leqslant 1} \| u (x) \|\tag{1}
\]

is finite. The mapping \(u \mapsto \| u \|\) is a norm on the vector space \(\mathcal{L}(\mathrm{E};\mathrm{F})\) of continuous linear mappings of \(\mathrm{E}\) in \(\mathrm{F}\) .

Let \(\mathbf{F}\) be a Banach space. Then \(\mathcal{L}(\mathrm{E};\mathrm{F})\) is a Banach space. The completion \(\hat{\mathrm{E}}\) of \(\mathrm{E}\) is a Banach space and the mapping \(u\mapsto u|\mathrm{E}\) is a bijective isometry of \(\mathcal{L}(\hat{\mathrm{E}};\mathrm{F})\) on \(\mathcal{L}(\mathrm{E};\mathrm{F})\) .

\begin{enumerate}
\def\labelenumi{\arabic{enumi})}
\setcounter{enumi}{1}
\tightlist
\item
  Let E be a normed space. Write \(E' = \mathcal{L}(E; K)\) where K carries the norm \(\lambda \mapsto |\lambda|\) . The Banach space \(E'\) is called the dual of E, and the dual \(E''\) of \(E'\) is called the bidual of E.
\end{enumerate}

Denote by \(\sigma(E, E')\) the coarsest topology on \(E\) for which all the linear forms \(x' \in E'\) are continuous; it is called the weakened topology of \(E\) . Denote by \(\sigma(E', E)\) the coarsest topology on \(E'\) for which the linear forms \(x' \mapsto \langle x', x \rangle\) on \(E'\) where \(x\) varies in \(E\) , are continuous; then \(\sigma(E', E)\) is called the weak topology on \(E'\) . The topology on \(E'\) deduced from the norm is called the strong topology.

\begin{enumerate}
\def\labelenumi{\arabic{enumi})}
\setcounter{enumi}{2}
\tightlist
\item
  Let E be a normed space and M be a closed vector subspace of E. Let \(\pi\) be the canonical mapping of E on E/M. A norm on the vector space E/M is defined by
\end{enumerate}

\[
\| \xi \| = \inf _ {\pi (x) = \xi} \| x \|.\tag{2}
\]

When E is a Banach space, then so also are M and E/M. For every normed space F, the linear mapping \(u \mapsto u \circ \pi\) of \(\mathcal{L}(\mathrm{E}/\mathrm{M}; \mathrm{F})\) in \(\mathcal{L}(\mathrm{E}; \mathrm{F})\) is isometric.

\begin{enumerate}
\def\labelenumi{\arabic{enumi})}
\setcounter{enumi}{3}
\tightlist
\item
  Let E be a normed space. For every \(x' \in E'\) , we have by definition
\end{enumerate}

\[
\| x^{\prime}\| = \sup_{\substack{\| x\| \leqslant 1\\ x\in \mathbf{E}}}\left|\langle x^{\prime},x\rangle \right|.\tag{3}
\]

Further (Hahn-Banach theorem), we have

\[
\| x \| = \sup_{\substack{\| x^{\prime}\| \leqslant 1\\ x^{\prime}\in \mathbf{E}^{\prime}}}\left|\langle x^{\prime},x\rangle \right|\tag{4}
\]

for all \(x \in E\) . In other words, the canonical mapping of E in its bidual \(E''\) is isometric.
